% !TeX root = ../../Book5.tex
% 纲要标题：## 第15章　泛包络代数与 PBW 定理
% 纲要位置：工作记录/计划纲要.md，第 4547 行。
\chapter{泛包络代数与 PBW 定理}
\label{b5:ch:15}
\BookChapterNumber{b5:ch:15}

\begin{chapterabstract}
本章构造 Lie 代数的泛包络代数，证明 PBW 定理及有序单项式基，随后研究伴随分次和 \(\mathfrak{sl}_2\) 的 Verma 模。
\end{chapterabstract}

\begin{learninggoals}
\begin{itemize}
\item 用泛性质构造包络代数同态，区分商代数、商模与左理想关系。
\item 分别证明 PBW 重排的终止性、重叠相容性与有序单项式的线性无关性。
\item 运用主符号计算维数、判断零因子并恢复 Lie 括号，区分 PBW 滤过与增广理想滤过。
\item 用右自由性证明诱导的正合性，以最高权向量刻画 Verma 模同态，并判断有限截断的相容性。
\end{itemize}
\end{learninggoals}

在 \(\mathfrak{sl}_2\) 中，生成元 \(e,f,h\) 满足 \(ef-fe=h\)。因此一个含 \(ef\) 的乘积可改写成含 \(fe\) 的乘积，却会多出一项次数更低的 \(h\)。若继续重排更长的字，可能有几条不同的重排次序；它们必须给出相同结果，才谈得上由有序单项式构成基。把 Lie 括号译成结合代数中的交换关系后，这个问题便成为泛包络代数的核心。PBW 定理不仅说每个元素能重排，还要证明重排后的表达不会暗含新的线性关系。

\ProseParagraphBreak
本章需要张量代数、滤过与伴随分次的基本定义，以及已经介绍的 Lie 代数表示和 \(\mathfrak{sl}_2\) 算子关系。PBW 主定理允许任意底域；只有对称化公式以及后面的复最高权例子另有特征零条件。

\ProseParagraphBreak
有序基与包络代数的经典证明可参阅 Humphreys\cite[Sections 17.3--17.4, pp. 91--94]{Humphreys1972LieAlgebras}。本章通过相邻重排及 Jacobi 恒等式处理重叠，直接证明 PBW 定理，再在所得基上计算 \(\mathfrak{sl}_2\) 的 Verma 模。这些构造也将用于后面的最高权分类。

% !TeX root = ../../Book5.tex
% 纲要标题：### 15.1 泛包络代数
% 纲要位置：工作记录/计划纲要.md，第 4549 行。

\section{泛包络代数}
\label{b5:sec:1501}


% 纲要范围：
%
% * 张量代数按 Lie 括号关系取商
% * 泛包络代数的泛性质
% * Lie 代数模与包络代数模
% * 用 Heisenberg 中心商比较多项式代数与 Weyl 代数；增广、余乘法和反自同构解释平凡模、张量模与对偶模
%

一个 Lie 代数表示只要求括号相容，而作用算子及其乘积属于自同态结合代数。泛包络代数以 Lie 代数为生成元，以 Lie 括号为乘法关系，使这些乘积可以在同一个代数中计算。取商构造之后，还须证明原 Lie 代数到泛包络代数的自然映射是单射；这将由 PBW 定理得到。

\subsection{张量代数按 Lie 括号关系取商}
\label{b5:subsec:1501:01}

设 \(\mathfrak g\) 为域 \(k\) 上 Lie 代数。张量代数
\[
	T(\mathfrak g)=\bigoplus_{r\ge0}\mathfrak g^{\otimes r}, \qquad \mathfrak g^{\otimes0}=k,
\]
以张量拼接为乘法，是由向量空间 \(\mathfrak g\) 自由生成的含幺结合代数。给定线性映射 \(\mathfrak g\to A\)，把纯张量送到各因子像的乘积，就唯一得到一个含幺代数同态 \(T(\mathfrak g)\to A\)。
\begin{definition}\label{b5:def:universal-enveloping}
	设 \(\mathfrak g\) 为域 \(k\) 上的 Lie 代数，记 \(T(\mathfrak g)\) 为其底层 \(k\) 向量空间的张量代数。令 \(I\) 为 \(T(\mathfrak g)\) 中由
	\[
		x\otimes y-y\otimes x-[x,y]\qquad(x,y\in\mathfrak g)
	\]
	生成的双边理想，其中 \([x,y]\) 看作次数一的张量。代数
	\[
		U(\mathfrak g)=T(\mathfrak g)/I
	\]
	称为 \(\mathfrak g\) 的\term[fanbaoluo-daishu]{泛包络代数}[universal enveloping algebra]。次数一部分的自然映射记为 \(\iota:\mathfrak g\to U(\mathfrak g)\)。
\end{definition}
\symidx{\(U(\mathfrak g)\)：Lie 代数的泛包络代数}
由关系立即得到 \(\iota\bigl([x,y]\bigr)=\iota(x)\iota(y)-\iota(y)\iota(x)\)， 所以 \(\iota\) 是到 \(U(\mathfrak g)^{-}\) 的 Lie 代数同态，其单射性将在 PBW 定理中证明。将全部正次数张量送到零、常数保持不变的映射零化 \(I\)，故下降为含幺代数同态 \(\varepsilon:U(\mathfrak g)\to k\)，称为增广映射。特别地，\(U(\mathfrak g)\) 的单位非零。

\ProseParagraphBreak

关系含有次数二与次数一的两部分，所以商一般不是由张量次数直接得到的分次代数。不过，“次数不超过 \(r\)”的子空间仍能传到商中，下一节将用这个滤过研究它。

\subsection{泛包络代数的泛性质}
\label{b5:subsec:1501:02}

\begin{theorem}\label{b5:thm:enveloping-universal}
	设 \(\mathfrak g\) 为域 \(k\) 上的 Lie 代数，\(A\) 为含幺结合 \(k\) 代数，\(A^{-}\) 为其交换子 Lie 代数，\(\iota:\mathfrak g\to U(\mathfrak g)\) 为泛包络代数的自然映射。对每个 Lie 代数同态 \(\phi:\mathfrak g\to A^{-}\)，存在唯一含幺 \(k\) 代数同态 \(\widetilde\phi:U(\mathfrak g)\to A\)，使 \(\widetilde\phi\iota=\phi\)。 也就是说，下图在底层 \(k\) 向量空间上交换，虚线额外要求为含幺代数同态，并在这一要求下唯一：
	\[
		\begin{tikzcd}[column sep=large,row sep=large]
			\mathfrak g\arrow[r,"\iota"]\arrow[dr,"\phi"'] &U(\mathfrak g)\arrow[d,dashed,"\widetilde\phi"]\\
			&A
		\end{tikzcd}
	\]
\end{theorem}
\begin{proof}
	先由张量代数的自由性延拓 \(\phi\)。因为 \(\phi\) 保持括号，每个关系元的像为 \(\phi(x)\phi(y)-\phi(y)\phi(x)-\phi\bigl([x,y]\bigr)=0\)。 所以延拓零化双边理想 \(I\)，下降到商。\(U(\mathfrak g)\) 由单位和 \(\iota(\mathfrak g)\) 生成，因此这些生成元上的值唯一决定代数同态。
\end{proof}
这就是泛包络代数的泛性质。它也使构造对同态自然：Lie 代数同态 \(f:\mathfrak g\to\mathfrak l\) 唯一诱导 \(U(f):U(\mathfrak g)\to U(\mathfrak l)\)，复合与恒等都因唯一性而保持。是否把单射映成单射，不能仅从泛性质推出；PBW 将另外给出这一结论。
\begin{proposition}\label{b5:prop:enveloping-quotient}
	设 \(\mathfrak g\) 为域 \(k\) 上的 Lie 代数，\(\mathfrak a\subseteq\mathfrak g\) 为理想，\(\iota:\mathfrak g\to U(\mathfrak g)\) 为泛包络代数的自然映射。令 \(J\) 为 \(U(\mathfrak g)\) 中由 \(\iota(\mathfrak a)\) 生成的双边理想，则
	\[
		U(\mathfrak g/\mathfrak a)\cong U(\mathfrak g)/J.
	\]
\end{proposition}
\begin{proof}
	映射 \(\mathfrak g\to U(\mathfrak g)/J\) 零化 \(\mathfrak a\)，故通过商 Lie 代数分解，再由泛性质延拓为左边到右边的同态。反过来，商映射 \(\mathfrak g\to\mathfrak g/\mathfrak a\) 诱导的包络代数同态零化 \(J\)，得到右边到左边的同态。两复合都在单位及来自 \(\mathfrak g\) 的生成元上为恒等，所以互逆。
\end{proof}
若 \(\mathfrak g\) 交换，定义中的关系只是让两个生成元交换，因而 \(U(\mathfrak g)=\operatorname{Sym}(\mathfrak g)\)。 特别地，一维交换 Lie 代数的包络代数是 \(k[t]\)，其表示可以由任意线性算子给出，不存在类似有限循环群关系 \(T^n=I\) 的限制。

\ProseParagraphBreak

取域 \(k\) 上的三维 Heisenberg Lie 代数 \(\mathfrak h=kx\oplus ky\oplus kz\)，规定 \([x,y]=z\)、\([x,z]=[y,z]=0\)。把基上的括号关系代入包络代数的定义，得到展示
\[
	U(\mathfrak h)\cong k\langle x,y,z\rangle/(xy-yx-z,\ xz-zx,\ yz-zy).
\]
这里各括号表示生成的双边理想，也用 \(x,y,z\) 表示它们在包络代数及其商中的自然像。中心 Lie 理想 \(kz\) 被置零后，\(\mathfrak h/kz\) 是由 \(x+kz,y+kz\) 生成的二维交换 Lie 代数。因此\cref{b5:prop:enveloping-quotient}在这个对象上给出
\[
	U(\mathfrak h)/(z)\cong U(\mathfrak h/kz)\cong k[X,Y].
\]
这个多项式识别可以直接核验：将商中的 \(x,y\) 送到 \(X,Y\)，三条展示关系都为零；反过来，将 \(X,Y\) 送到商中的 \(x,y\)，因两者交换而得到多项式代数同态。两个复合固定单位和各生成元，故互逆。将括号的中心值置零，恰好使原来的两个不交换生成元成为多项式变量。

\ProseParagraphBreak

若改在结合代数中施加 \(z=1\)，则得到第一 Weyl 代数
\[
	U(\mathfrak h)/(z-1)\cong A_1(k)\coloneqq k\langle p,q\rangle/(pq-qp-1).
\]
从左边到右边的同态由 \(x\mapsto p\)、\(y\mapsto q\)、\(z\mapsto1\) 给出；反向同态由 \(p\mapsto x\)、\(q\mapsto y\) 给出，因为左边满足 \(xy-yx=1\)。两复合仍固定生成元，故互逆。在 \(k[t]\) 上令
\[
	D(P)=\frac{\mathrm{d}P}{\mathrm{d}t},\qquad T(P)=tP.
\]
对任意 \(P\in k[t]\)，乘积求导给出 \(D(T(P))-T(D(P))=P\)。所以 \(x\mapsto D\)、\(y\mapsto T\)、\(z\mapsto I\) 给出 \(U(\mathfrak h)/(z-1)\) 在 \(k[t]\) 上的幺性左模；其交换子为非零的恒等算子，说明这个商仍非交换。这里指定的是中心元素的标量作用。由 \(\varepsilon(z-1)=-1\ne0\)，理想 \((z-1)\) 不可能由某个 Lie 理想的自然像生成，因为所有这类自然像及其生成的双边理想都被 \(\varepsilon\) 零化。

\subsection{Lie 代数模与包络代数模}
\label{b5:subsec:1501:03}

\begin{theorem}\label{b5:thm:lie-enveloping-modules}
	设 \(\mathfrak g\) 为域 \(k\) 上的 Lie 代数，\(\iota:\mathfrak g\to U(\mathfrak g)\) 为泛包络代数的自然映射。左 \(\mathfrak g\) 模与幺性左 \(U(\mathfrak g)\) 模给出同一个范畴：在相同 \(k\) 向量空间上，作用通过 \(\iota\) 限制或按泛性质唯一延拓；同态、子模、商模与正合列均保持。
\end{theorem}
\begin{proof}
	一个 Lie 代数表示 \(\rho:\mathfrak g\to\mathfrak{gl}(V)\) 是到结合代数 \(\operatorname{End}_k(V)\) 的交换子 Lie 代数的同态，故由\cref{b5:thm:enveloping-universal}唯一延拓为 \(U(\mathfrak g)\) 的含幺作用。反向限制显然恢复表示条件，唯一性使两种过程互逆。 设 \(V,W\) 的 Lie 代数作用分别为 \(\rho_V,\rho_W\)。一个 \(k\) 线性映射 \(T:V\to W\) 是 Lie 代数模同态，当且仅当
	\[
		T\rho_V(x)=\rho_W(x)T\qquad(x\in\mathfrak g).
	\]
	反复使用此式，得到
	\[
		T\rho_V(x_1)\cdots\rho_V(x_r) =\rho_W(x_1)\cdots\rho_W(x_r)T.
	\]
	这些乘积和单位的线性组合给出整个 \(U(\mathfrak g)\) 的作用，所以 \(T\) 也与延拓后的作用相容。反向限制则立即成立。因此两种同态相同，子空间的不变性也相同，核、商及正合性因而不变。
\end{proof}
包络代数中的乘积由连续施行 Lie 代数算子得到。例如在 \(\mathfrak{sl}_2\) 模上，\(ef\)、\(h^2\) 或它们的和都是自然的算子，但一般不是某个 Lie 代数元素的作用。上一章的 Casimir 算子正需要这种乘积；现在可以把它看成一个与具体表示无关的代数元素。

\ProseParagraphBreak

两个包络代数左模的张量积取在域 \(k\) 上。要在这个空间上作用一个代数元素，需要把它分配到两个因子；要在对偶空间上作用，则需要倒转乘积次序。这两项构造和增广映射都可以由同一个泛性质得到。
\begin{proposition}\label{b5:prop:enveloping-tensor-antipode}
	设 \(\mathfrak g\) 为域 \(k\) 上的 Lie 代数，\(U=U(\mathfrak g)\)，\(\iota:\mathfrak g\to U\) 为自然映射。存在唯一含幺 \(k\) 代数同态
	\[
		\varepsilon:U\longrightarrow k,\qquad \Delta:U\longrightarrow U\otimes_k U,
	\]
	以及唯一 \(k\) 线性反自同构 \(S:U\to U\)，满足
	\[
		\begin{aligned}
			\varepsilon(\iota(x))&=0,\\
			\Delta(\iota(x))&=\iota(x)\otimes1+1\otimes\iota(x),\\
			S(\iota(x))&=-\iota(x) \qquad(x\in\mathfrak g).
		\end{aligned}
	\]
	其中 \(S(1)=1\)、\(S(ab)=S(b)S(a)\)。这些映射满足
	\[
		\begin{gathered}
			(\Delta\otimes\operatorname{id}_U)\Delta =(\operatorname{id}_U\otimes\Delta)\Delta,\qquad S^2=\operatorname{id}_U,\\
			(\varepsilon\otimes\operatorname{id}_U)\Delta =\operatorname{id}_U =(\operatorname{id}_U\otimes\varepsilon)\Delta,
		\end{gathered}
	\]
	最后一行采用 \(k\otimes_k U\cong U\cong U\otimes_k k\) 的自然识别。

	若 \(V,W\) 为幺性左 \(U\) 模，则下列公式分别给出 \(k\)、\(V\otimes_k W\)、\(V^*=\operatorname{Hom}_k(V,k)\) 上的幺性左 \(U\) 模结构：
	\[
		\begin{aligned}
			u\cdot c&=\varepsilon(u)c,\\
			u\cdot(v\otimes w)&=\sum_r a_rv\otimes b_rw \quad\left(\Delta(u)=\sum_r a_r\otimes b_r\right),\\
			(u\cdot\varphi)(v)&=\varphi\bigl(S(u)v\bigr).
		\end{aligned}
	\]
	这里 \(u\in U\)、\(c\in k\)、\(v\in V\)、\(w\in W\)、\(\varphi\in V^*\)；张量的求和为有限和。
\end{proposition}
\symidx{\(\varepsilon,\Delta,S\)：包络代数的增广、余乘法与反自同构}
\begin{proof}
	零 Lie 代数同态 \(\mathfrak g\to k^{-}\) 由\cref{b5:thm:enveloping-universal}唯一延拓为 \(\varepsilon\)，它与前面构造的增广映射相同。对 \(\Delta\)，在 \(U\otimes_k U\) 中，不同因子的元素交换，所以
	\[
		\begin{aligned}
			&[\iota(x)\otimes1+1\otimes\iota(x),\, \iota(y)\otimes1+1\otimes\iota(y)]\\
			&\qquad=\iota([x,y])\otimes1+1\otimes\iota([x,y]).
		\end{aligned}
	\]
	因此指定的线性映射保持 Lie 括号，同一泛性质给出良定义且唯一的 \(\Delta\)。两个三重余乘法都是到 \(U^{\otimes3}\) 的含幺代数同态，并将 \(\iota(x)\) 送到
	\[
		\iota(x)\otimes1\otimes1 +1\otimes\iota(x)\otimes1 +1\otimes1\otimes\iota(x).
	\]
	它们在生成元和单位上一致，故相等。用 \(c\otimes u\mapsto cu\)、\(u\otimes c\mapsto cu\) 作自然识别后，两种含 \(\varepsilon\) 的复合都固定 \(\iota(x)\) 和单位，同样得到余单位等式。

	对 \(S\)，把目标换为反代数 \(U^{\mathrm{op}}\)，其乘法次序与 \(U\) 相反。于是
	\[
		[-\iota(x),-\iota(y)]_{\mathrm{op}} =\iota(y)\iota(x)-\iota(x)\iota(y) =-\iota([x,y]).
	\]
	所以 \(x\mapsto-\iota(x)\) 是到 \((U^{\mathrm{op}})^{-}\) 的 Lie 代数同态，由泛性质唯一延拓为 \(U\to U^{\mathrm{op}}\) 的含幺代数同态，也就是反同态 \(S\)。它在生成元乘积上满足
	\[
		S\bigl(\iota(x_1)\cdots\iota(x_r)\bigr) =(-1)^r\iota(x_r)\cdots\iota(x_1).
	\]
	\(S^2\) 为代数同态，并固定所有生成元和单位，故为恒等。这样 \(S\) 可逆，确为反自同构。

	\(\varepsilon\) 是含幺代数同态，所以第一条作用给出左模。映射
	\[
		U\otimes_k U\longrightarrow\operatorname{End}_k(V\otimes_k W),\qquad a\otimes b\longmapsto\bigl(v\otimes w\mapsto av\otimes bw\bigr)
	\]
	由双线性得到，并保持乘法和单位。它与 \(\Delta\) 复合，给出第二条作用；这也说明该公式不依赖 \(\Delta(u)\) 的具体张量分解。对第三条作用，若 \(u,a\in U\)，则
	\[
		\begin{aligned}
			\bigl(u\cdot(a\cdot\varphi)\bigr)(v) &=\varphi\bigl(S(a)S(u)v\bigr)\\
			&=\varphi\bigl(S(ua)v\bigr) =\bigl((ua)\cdot\varphi\bigr)(v).
		\end{aligned}
	\]
	单位同样作用为恒等，故对偶也得到幺性左模。
\end{proof}
限制到 \(\mathfrak g\) 后，第一条作用使每个 \(x\) 为零，给出平凡模；后两条则恢复\cref{b5:prop:lie-dual-tensor-hom}中的张量和对偶作用。以下把 \(\iota(x)\) 简写为 \(x\)。因为 \(\Delta\) 保持乘法，
\[
	\Delta(x^2)=x^2\otimes1+2x\otimes x+1\otimes x^2,
\]
所以
\[
	x^2(v\otimes w)=x^2v\otimes w+2xv\otimes xw+v\otimes x^2w.
\]
特征二时中间项为零；任意特征下，高次乘积的张量作用都由生成元的作用和 \(\Delta\) 的乘法相容性确定。

\ProseParagraphBreak

对偶上的负号来自 \(S(x)=-x\)：\((x\varphi)(v)=-\varphi(xv)\)，而反同态性质保证继续作用时乘积次序正确。例如 \((xy\varphi)(v)=\varphi(yxv)\)。这个负号也与一次项的取逆相容：在 \(U\otimes_k k[\delta]/(\delta^2)\) 中，\((1+\delta x)^{-1}=1-\delta x\)。增广 \(\varepsilon\)、余乘法 \(\Delta\) 和反自同构 \(S\) 因而分别控制平凡模、张量模与对偶模；\chapref{b5:ch:21}将把它们作为同一套 Hopf 代数结构数据统一讨论。

% !TeX root = ../../Book5.tex
% 纲要标题：### 15.2 滤过与 PBW
% * 新增小节状态：初稿已完成；任意特征的标准辛括号、Hamilton导子及中心计算均就地展开，新增4道习题及完整解答；新增范围尚未正式审核，原有内容保留。
% 纲要位置：工作记录/计划纲要.md，第 4556 行。

\section{滤过与 PBW}
\label{b5:sec:1502}


% 纲要范围：
%
% * 标准滤过
% * Poincaré–Birkhoff–Witt 定理
% * 伴随分次与对称代数
% * 正文选择有序单项式与滤过的独立 PBW 证明，所需重排与歧义相容性论证就地完成；第二卷第20.6节的 Diamond 引理仅提供另一证明，不是必需的选读内容
% * 用实际重叠与 Jacobi 失败比较正规式；主符号的 Poisson 括号保留 Lie 括号，增广一阶商与交换化及 PBW 滤过分开
%

张量次数在关系中可以下降，却不会上升。把这种现象写成滤过后，原来的非交换乘法在最高次上变成交换乘法。PBW 定理进一步断言：在最高次上，除了交换关系以外没有新的关系。

\subsection{标准滤过}
\label{b5:subsec:1502:01}

令 \(F_dU(\mathfrak g)\) 为长度不超过 \(d\) 的张量在商中的像，\(d\ge0\)，并令 \(F_{-1}=0\)。则
\[
	F_0=k,\qquad F_dF_e\subseteq F_{d+e},\qquad U(\mathfrak g)=\bigcup_{d\ge0}F_d.
\]
这称为标准滤过。伴随分次代数为
\[
	\operatorname{gr}U(\mathfrak g) =\bigoplus_{d\ge0}F_d/F_{d-1}.
\]
乘法由滤过给出：对 \(u\in F_d\)、\(v\in F_e\)，规定
\[
	(u+F_{d-1})(v+F_{e-1})=uv+F_{d+e-1}.
\]
更换代表产生的误差分别落在 \(F_{d-1}F_e\) 与 \(F_dF_{e-1}\)，都包含于 \(F_{d+e-1}\)，所以这一乘法良定义。 由基本关系 \(\iota(x)\iota(y)-\iota(y)\iota(x)=\iota\bigl([x,y]\bigr)\in F_1\)，两个次数一的主符号交换。更一般地，利用 \([ab,c]=a[b,c]+[a,c]b\) 逐次展开，可得
\[
	[F_d,F_e]\subseteq F_{d+e-1}.
\]
因此伴随分次代数交换，次数一的符号给出满同态
\[
	\operatorname{Sym}(\mathfrak g)\longrightarrow\operatorname{gr}U(\mathfrak g).
\]
满射只说明最高次的生成性；它是否单射，正是下面要证明的内容。

\subsection{Poincaré–Birkhoff–Witt 定理}
\label{b5:subsec:1502:02}

这个定理以 Poincaré、Birkhoff 与 Witt 命名。\footnote{庞加莱（Henri Poincaré，1854-1912）是法国数学家，研究拓扑、动力系统及数学物理；加勒特·伯克霍夫（Garrett Birkhoff，1911-1996）与维特（Ernst Witt，1911-1991）分别是美国和德国数学家，研究代数结构。}
\begin{theorem}[Poincaré--Birkhoff--Witt]\label{b5:thm:pbw}
	设 \(\mathfrak g\) 为任意域 \(k\) 上的 Lie 代数，并选定一组带全序的基 \((x_i)_{i\in I}\)，其中 \(I\) 为全序指标集。记 \(\iota:\mathfrak g\to U(\mathfrak g)\) 为自然映射，并按乘积长度给 \(U(\mathfrak g)\) 取张量次数滤过。在 \(U(\mathfrak g)\) 中，单位与全部有序单项式
	\[
		\iota(x_{i_1})\cdots\iota(x_{i_r}), \qquad r\ge1,\quad i_1\le\cdots\le i_r,
	\]
	构成一组 \(k\) 基。因此自然映射 \(\operatorname{Sym}(\mathfrak g)\to\operatorname{gr}U(\mathfrak g)\) 是分次代数同构，\(\iota:\mathfrak g\to U(\mathfrak g)\) 单射。
\end{theorem}
\symidx{\(\operatorname{gr}U(\mathfrak g)\)：泛包络代数的 PBW 伴随分次代数}
\begin{proof}
	先令 \(\mathcal O\) 为以全部有序字（包括空字）为形式基的自由 \(k\) 向量空间。把每个形式字送到相应张量字，得到线性映射 \(j:\mathcal O\to T(\mathfrak g)\)；张量字是一组基，所以这时 \(j\) 是单射。我们尚未把 \(\mathcal O\) 看作 \(U(\mathfrak g)\) 的子空间，而要先构造 \(\operatorname{NF}:T(\mathfrak g)\to\mathcal O\)，再证明它能够下降到商。 在张量代数的字上，遇到逆序的相邻因子 \(x_jx_i\)、\(j>i\)，便使用重排
	\[
		x_jx_i\ \longmapsto\ x_ix_j+[x_j,x_i],
	\]
	把右端括号按所选基展开。对一个字，按字典序比较长度与逆序对数。交换项保持长度而减少一个逆序对，括号项长度减少一，因此每一步产生的字都比原字小。若起始字的长度为 \(N\)，则后续字的长度 \(\ell\le N\)，逆序对数至多为 \(\ell(\ell-1)/2\)。这些复杂度只有有限种可能，故每条重写分支的步数都有统一上界。每次括号展开只有有限个非零基系数，整棵重写树因而既有限分支又有界深度，所以最终得到有限项的线性组合。 下面证明结果与选择哪一对先重排无关。对长度及逆序数作归纳。两次重排若发生在不相交的位置，分别展开两处会给出相同的四组项；归纳假设保证其余较小项继续重排的结果相同。若两处相交，就只能是三个相邻基元 \(x_kx_jx_i\) 的左右两个相邻对；两处都能重排要求 \(k>j\) 且 \(j>i\)，所以唯一的重叠情形是 \(k>j>i\)。先处理左边再继续交换，得到
	\[
		x_ix_jx_k+[x_j,x_i]x_k+x_j[x_k,x_i]+[x_k,x_j]x_i;
	\]
	先处理右边，则得到
	\[
		x_ix_jx_k+x_i[x_k,x_j]+[x_k,x_i]x_j+x_k[x_j,x_i].
	\]
	两式之差都是长度二的项。在归纳中，任意 \(u,v\in\mathfrak g\)（视为张量代数中的一次元素）已满足 \(\operatorname{NF}(uv-vu)=\operatorname{NF}\bigl([u,v]\bigr)\)：先对基元检查所选次序，再双线性延拓即可。因此这个差的正规式等于
	\[
		\operatorname{NF}\Bigl( \bigl[[x_j,x_i],x_k\bigr]+\bigl[x_j,[x_k,x_i]\bigr]+\bigl[[x_k,x_j],x_i\bigr] \Bigr)=0
	\]
	其中括号内的和由 Jacobi 恒等式为零。若前后还有其他因子，同样的差只含比原字短一的项，故在相同归纳范围内。这证明所有重写选择相容，给出到有序字空间的线性正规式映射 \(\operatorname{NF}\)。 每次重排之差都属于定义理想 \(I\)，所以对任意 \(u\in T(\mathfrak g)\)，有 \(u-j(\operatorname{NF}(u))\in I\)。反过来，对于任意字 \(u,v\) 及基元 \(x_i,x_j\)，
	\[
		\operatorname{NF}\Bigl(u\bigl(x_jx_i-x_ix_j-[x_j,x_i]\bigr)v\Bigr)=0,
	\]
	当 \(j>i\) 时，可以先在中间重排这一对因子，再正规化；前面证明的选择无关性给出
	\[
		\operatorname{NF}(ux_jx_iv) =\operatorname{NF}\bigl(u(x_ix_j+[x_j,x_i])v\bigr).
	\]
	移项即得上述等式。\(j<i\) 时交换两个指标，\(j=i\) 时定义关系本身为零，结论仍成立。由双线性及理想生成性，\(I\subseteq\ker\operatorname{NF}\)，于是 \(\operatorname{NF}\) 才下降为 \(\overline{\operatorname{NF}}:U(\mathfrak g)\to\mathcal O\)。记商映射为 \(q:T(\mathfrak g)\to U(\mathfrak g)\)。有序字被原样保留，所以 \(\operatorname{NF}j=1_{\mathcal O}\)，即 \(\overline{\operatorname{NF}}\,qj=1_{\mathcal O}\)；前面的模 \(I\) 等式则给出 \(qj\,\overline{\operatorname{NF}}=1_{U(\mathfrak g)}\)。这两个线性映射互逆，才证明有序字的像在商中生成且线性无关。 按长度分组，长度恰为 \(d\) 的这些基元给出 \(F_d/F_{d-1}\) 的基，与对称代数的单项式基对应，故自然满射为同构。长度一的基元因此仍线性无关，得到 \(\iota\) 单射。
\end{proof}
这称为 PBW 定理。证明只使用交替性、Jacobi 及重排，没有除以整数，所以不需要特征零。以下在不会混淆时省去 \(\iota\)，把 \(\mathfrak g\) 作为 \(U(\mathfrak g)\) 的线性子空间。

\ProseParagraphBreak

可以在 \(\mathbb C\) 上的 \(\mathfrak{sl}_2\) 中直接看一次重叠。取基次序 \(f<h<e\)，由
\[
	eh=(h-2)e,\qquad hf=f(h-2),\qquad ef=fe+h
\]
重排 \(ehf\)。先处理 \(eh\)，得到
\[
	ehf=(h-2)ef=(h-2)(fe+h) =f(h-4)e+h(h-2).
\]
先处理 \(hf\)，则得到
\[
	ehf=ef(h-2)=(fe+h)(h-2) =f(h-4)e+h(h-2).
\]
两条路径的结果相同，正是证明中三重重叠相容性的一个具体实例。

\ProseParagraphBreak

有序单项式基依赖于所选的 Lie 代数基及其次序，PBW 的分次同构却不依赖这些选择：它由自然映射 \(\mathfrak g\to F_1/F_0\) 诱导。例如对满足 \([h,e]=e\) 的二维 Lie 代数，次序 \(e<h\) 给出基 \(e^ah^b\)，次序 \(h<e\) 给出基 \(h^be^a\)。在包络代数中 \(he=eh+e\)，两组基并不相同；取次数二的符号时，低次项 \(e\) 消失，两者都对应 \(\operatorname{Sym}^2(\mathfrak g)\) 中的同一个单项式。这里忽略的是当前次数以下的部分，并没有把原包络代数变成交换代数。

\begin{lemma}\label{b5:lem:failed-jacobi-collapse}
	设 \(k\) 为任意域，结合代数
	\[
		A=k\langle x,y,z\rangle/ (xy-yx-x,\ yz-zy-y,\ zx-xz-z).
	\]
	则三个生成元在 \(A\) 中的像都为零，因而 \(A\cong k\)。
\end{lemma}
\begin{proof}
	仍用 \(x,y,z\) 表示商中的像，括号表示结合代数的交换子。结合代数的交换子必满足 Jacobi 恒等式，而给定关系使
	\[
		[x,[y,z]]+[y,[z,x]]+[z,[x,y]]=x+y+z.
	\]
	所以 \(z=-x-y\)。代入 \([y,z]=y\)，得到 \([y,-x-y]=x=y\)；再由 \([x,y]=x\) 得到 \(x=0\)，故 \(y=z=0\)。这只用了加法与交换子，没有除以任何整数。于是 \(A\) 由单位生成，而把三个生成元送到零的增广同态 \(A\to k\) 保证单位仍非零，得到所述同构。
\end{proof}
这个例子说明 Jacobi 条件在 PBW 中保护的正是生成元的线性无关性。若任意指定一次交换子关系而不满足 Jacobi，重排可能迫使生成元本身消失。

\subsection{伴随分次与对称代数}
\label{b5:subsec:1502:03}

\begin{corollary}\label{b5:cor:pbw-dimension-domain}
	设 \(\mathfrak g\) 为域 \(k\) 上的 \(n\) 维 Lie 代数，\(F_dU(\mathfrak g)\) 为其泛包络代数中由长度不超过整数 \(d\ge0\) 的乘积生成的 \(k\) 子空间，其中长度零的乘积是单位。则
	\[
		\dim F_dU(\mathfrak g)=\binom{n+d}{n}\qquad(d\ge0),
	\]
	且 \(U(\mathfrak g)\) 没有零因子。若 \(\mathfrak g\ne0\)，则 \(U(\mathfrak g)\) 无限维。
\end{corollary}
\begin{proof}
	PBW 基把问题化为计算 \(n\) 个变量中全次数不超过 \(d\) 的单项式数；加一个记录剩余次数的变量后，用非负整数解计数得到 \(\binom{n+d}{n}\)。\(n=0\) 时两侧均为一。 若 \(u,v\ne0\)，取它们在标准滤过中的最高次主符号。由 PBW，伴随分次是多项式整环，两个非零主符号的乘积不为零，所以 \(uv\ne0\)。若 \(n\ge1\)，某个基元的任意不同次幂都是 PBW 基元，故代数无限维。
\end{proof}
\begin{proposition}\label{b5:prop:pbw-poisson}
	设 \(\mathfrak g\) 为任意域 \(k\) 上的 Lie 代数，\(F_d\) 为 \(U(\mathfrak g)\) 的标准滤过，记 \(\sigma_d(u)=u+F_{d-1}\)。对非负整数 \(r,s\) 及 \(u\in F_r\)、\(v\in F_s\)，当 \(r+s\ge1\) 时，公式
	\[
		\{\sigma_r(u),\sigma_s(v)\} =\sigma_{r+s-1}(uv-vu)
	\]
	良定义；两个次数零元素的括号规定为零，再按双线性延拓到 \(\operatorname{gr}U(\mathfrak g)\)。这个括号交替、满足 Jacobi 恒等式，且
	\[
		\{a,bc\}=\{a,b\}c+b\{a,c\}.
	\]
	因此它是次数 \(-1\) 的 Poisson 括号。在 PBW 同构 \(\operatorname{gr}U(\mathfrak g)\cong\operatorname{Sym}(\mathfrak g)\) 下，对 \(x,y\in\mathfrak g\) 有 \(\{x,y\}=[x,y]\)。
\end{proposition}
具有上述性质的括号称为\term[Poisson-kuohao]{Poisson 括号}[Poisson bracket]，它把交换乘法与 Lie 括号通过导子规则联系起来。\footnote{泊松（Siméon-Denis Poisson，1781-1840），法国数学家、物理学家，研究分析、概率及力学；力学中的泊松括号以他命名。}
\symidx{\(\{a,b\}\)：\(\operatorname{Sym}(\mathfrak g)\) 上由 Lie 括号诱导的 Poisson 括号}
\begin{proof}
	已经证明 \([F_r,F_s]\subseteq F_{r+s-1}\)。若把 \(u\) 换成 \(u+u'\)，其中 \(u'\in F_{r-1}\)，交换子的变化属于 \(F_{r+s-2}\)，在所取商中为零；更换 \(v\) 同理。次数零的元素是标量，与所有元素交换，所以涉及标量的括号也为零。 对齐次元素，交替性与反对称性来自 \([u,u]=0\) 及 \([u,v]=-[v,u]\)。对非齐次元素 \(a=\sum_r a_r\)，\(\{a,a\}\) 中同次项为零，不同次项两两相消，故仍为零，包括特征二的情形。交换子的 Jacobi 恒等式及
	\[
		[u,vw]=[u,v]w+v[u,w]
	\]
	在相应次数上取符号，分别给出 Jacobi 恒等式与导子规则。最后，对次数一的 \(x,y\)，包络代数的定义关系给出 \(xy-yx=[x,y]\in F_1\)，其次数一符号就是原 Lie 括号。
\end{proof}
交换子的二次最高项在伴随分次中消失，但低一度的首个误差由这个括号保留下来。例如 Heisenberg 代数满足 \([x,y]=z\)、\([x,z]=[y,z]=0\)，于是其对称代数 \(k[x,y,z]\) 中
\[
	\{x,y\}=z,\qquad \{x,z\}=\{y,z\}=0, \qquad \{x^2,y\}=2xz.
\]
最后一个公式来自导子规则；在特征二中它的右端为零。

\begin{proposition}\label{b5:prop:pbw-symmetrization}
	设 \(k\) 为特征零域，\(\mathfrak g\) 为 \(k\) 上的 Lie 代数，\(\operatorname{Sym}(\mathfrak g)\) 为其对称代数。对整数 \(r\ge0\) 及 \(x_1,\ldots,x_r\in\mathfrak g\)，在右端通过自然映射把各 \(x_i\) 看作 \(U(\mathfrak g)\) 中的元素，则对称化公式
	\[
		\operatorname{sym}(x_1\cdots x_r) =\frac1{r!}\sum_{\sigma\in S_r}x_{\sigma(1)}\cdots x_{\sigma(r)}
	\]
	给出 \(\operatorname{Sym}(\mathfrak g)\to U(\mathfrak g)\) 的滤过向量空间同构。它与伴随作用相容，但一般不是代数同构。
\end{proposition}
\symidx{\(\operatorname{sym}\)：对称代数到泛包络代数的 PBW 对称化映射}
\begin{proof}
	右端对 \(x_1,\ldots,x_r\) 对称且多线性，故由对称幂的泛性质良定义。给 \(\operatorname{Sym}(\mathfrak g)\) 取全次数滤过；对称化的最高次符号等于 \(x_1\cdots x_r\)，所以诱导的分次映射在 PBW 同构下为恒等。 对任意 \(u\in F_dU(\mathfrak g)\)，由这个分次同构可取 \(s\in\operatorname{Sym}^d(\mathfrak g)\)，使 \(u-\operatorname{sym}(s)\in F_{d-1}\)。从 \(F_{-1}=0\) 开始对 \(d\) 归纳，便可逐次匹配各次主符号，得到满射。若 \(s\ne0\)，取其最高非零齐次次数 \(d\)，则其主符号非零；分次映射的单射性保证 \(\operatorname{sym}(s)\) 的 \(d\) 次主符号也非零，故 \(\operatorname{sym}(s)\ne0\)。这证明了单射性。 伴随作用在张量乘积上按导子规则逐项作用，并与排列因子交换，所以对称化保持作用。其乘法差异在次数二已经出现：
	\[
		\operatorname{sym}(xy)=\frac{xy+yx}{2} =xy-\frac{[x,y]}2, \qquad \operatorname{sym}(x)\operatorname{sym}(y)=xy.
	\]
	若原括号不为零，PBW 的单射性保证这个差不为零，因此对称化不是代数同态。
\end{proof}
在 Heisenberg 例子中，差正是 \(-z/2\)。它低于 \(xy\) 的次数，因而在伴随分次中消失；这解释了对称化为何能诱导分次代数同构，却通常只给出原代数之间的向量空间同构。

\subsection{有序单项式与重排方法}
\label{b5:subsec:1502:04}

对于以所选基表示的具体元素，依次重排即可得到 PBW 系数；定理保证结果与重排路径无关。也可对照第二卷第20.6节的非交换重写与 Diamond 引理，按重写歧义组织这里的重叠检查。

\ProseParagraphBreak

乘积长度与增广理想的幂却不能混为一谈。记 \(\varepsilon:U(\mathfrak g)\to k\) 为把 \(\mathfrak g\) 送到零的增广同态，\(I=\ker\varepsilon\)。长度一的元素 \([x,y]\) 同时等于两个长度二乘积之差 \(xy-yx\)，所以它在 \(I/I^2\) 中为零，在 \(F_1/F_0\) 中却未必为零。

\begin{proposition}\label{b5:prop:enveloping-augmentation-cotangent}
	设 \(\mathfrak g\) 为域 \(k\) 上的 Lie 代数，\(I\) 为 \(U(\mathfrak g)\) 的增广理想。自然映射 \(x\mapsto x+I^2\) 诱导向量空间同构
	\[
		\mathfrak g/[\mathfrak g,\mathfrak g]\xrightarrow{\sim} I/I^2.
	\]
	此外，保持单位的 \(k\) 代数同态 \(U(\mathfrak g)\to k\)，也即 \(\mathfrak g\) 的一维表示，与 \(\bigl(\mathfrak g/[\mathfrak g,\mathfrak g]\bigr)^*\) 的元素一一对应。 若 \(\mathfrak g\ne0\) 且 \([\mathfrak g,\mathfrak g]=\mathfrak g\)，则 \(I=I^2\ne0\)，并且 \(\bigcap_{r\ge1}I^r=I\ne0\)。
\end{proposition}
\symidx{\(I/I^2\)：泛包络代数增广理想的商，与 Lie 代数的交换化同构}
\begin{proof}
	增广理想由 \(\mathfrak g\) 生成，也由所有正长度的乘积线性张成。模掉 \(I^2\) 后，长度至少二的乘积为零，所以 \(\mathfrak g\to I/I^2\) 满射。关系 \([x,y]=xy-yx\in I^2\) 又说明它通过交换化 \(\mathfrak g_{\mathrm{ab}}=\mathfrak g/[\mathfrak g,\mathfrak g]\)。 为检查没有更多一次关系，取平方零代数 \(B=k\oplus\mathfrak g_{\mathrm{ab}}\)，乘法为
	\[
		(a,v)(b,w)=(ab,aw+bv).
	\]
	映射 \(x\mapsto(0,\bar x)\) 是从 \(\mathfrak g\) 到 \(B^-\) 的 Lie 同态，故延拓为代数同态 \(U(\mathfrak g)\to B\)。它把 \(I\) 送入 \(0\oplus\mathfrak g_{\mathrm{ab}}\)，并杀掉 \(I^2\)。诱导的 \(I/I^2\to\mathfrak g_{\mathrm{ab}}\) 把 \(x+I^2\) 送回 \(\bar x\)，是前面满射的逆映射，得到同构。 一个 \(k\) 代数同态 \(U(\mathfrak g)\to k\) 限制到 \(\mathfrak g\) 后是线性泛函 \(\ell\)，且必须满足 \(\ell([x,y])=0\)，因为目标交换。反过来，任何满足这个条件的 \(\ell\) 都是到交换 Lie 代数 \(k\) 的 Lie 同态，由包络代数的泛性质唯一延拓。因此这些同态恰由交换化的对偶参数化。 若 \([\mathfrak g,\mathfrak g]=\mathfrak g\)，则已经证明的同构给出 \(I/I^2=0\)，即 \(I=I^2\)。归纳得 \(I^r=I\) 对所有 \(r\ge1\) 成立，而当 \(\mathfrak g\ne0\) 时，PBW 的单射性保证 \(I\ne0\)，得到最后的结论。
\end{proof}
同一个商也能记录增广附近的一阶变化。令 \(D=k[\delta]/(\delta^2)\)；保持单位、模 \(\delta\) 后等于 \(\varepsilon\) 的代数同态 \(U(\mathfrak g)\to D\)，在生成元上必为 \(x\mapsto\delta\ell(x)\)。由于 \(\delta^2=0\)，定义关系要求 \(\ell([x,y])=0\)，所以这里的 \(\ell\) 仍然来自 \(\mathfrak g_{\mathrm{ab}}^*\)。

\ProseParagraphBreak

增广理想的幂因此可能无法逐层分辨非零元素。例如在 \(\mathbb C\) 上的 \(\mathfrak{sl}_2\) 中，
\[
	h=ef-fe,\qquad e=\frac{he-eh}{2},\qquad f=-\frac{hf-fh}{2}
\]
使三个生成元都属于 \(I^2\)。与此同时，标准滤过仍有 \(F_1/F_0\cong\mathfrak{sl}_2\)，维数为三；它按乘积长度保留的一次符号，并不是 \(I/I^2\) 所记录的一次信息。

% 写作状态：15.2.5正文、证明、算例及配套习题已完成初稿，尚未正式审核。
% 写作范围：复用本节已有PBW Poisson括号，补Poisson代数、标准辛括号及其计算。
% #### 15.2.1 标准滤过
% #### 15.2.2 Poincaré–Birkhoff–Witt 定理
% #### 15.2.3 伴随分次与对称代数
% #### 15.2.4 有序单项式与重排方法
% #### 15.2.5 Poisson 代数与标准辛括号

\subsection{Poisson 代数与标准辛括号}
\label{b5:subsec:1502:05}

\cref{b5:prop:pbw-poisson}在对称代数上保留了原来的 Lie 括号。这个构造包含两种相容的运算：多项式照常相乘，而每个元素又通过括号给出一个导子。把这两种运算单独列出，便能讨论不一定来自某个包络代数标准滤过的例子。

\begin{definition}\label{b5:def:poisson-algebra}
	设 \(k\) 为域，\(A\) 为交换含幺 \(k\) 代数。若 \(k\) 双线性运算 \(\{-,-\}:A\times A\to A\) 对任意 \(a,b,c\in A\) 满足
	\[
		\begin{gathered}
			\{a,a\}=0,\qquad
			\{a,\{b,c\}\}+\{b,\{c,a\}\}+\{c,\{a,b\}\}=0,\\
			\{a,bc\}=\{a,b\}c+b\{a,c\},
		\end{gathered}
	\]
	就称带有这个括号的 \(A\) 为\term[Poisson-daishu]{Poisson 代数}[Poisson algebra]。两个 Poisson 代数之间保持括号的含幺 \(k\) 代数同态称为 Poisson 同态。若理想 \(I\subseteq A\) 还满足 \(\{A,I\}\subseteq I\)，则称它为\term[Poisson-lixiang]{Poisson 理想}[Poisson ideal]。
\end{definition}

交替性作用于 \(a+b\) 后给出 \(\{a,b\}=-\{b,a\}\)，所以括号在第一变量也满足乘积法则。特征二中反对称性不能推出交替性，定义仍要求 \(\{a,a\}=0\)。由 \(\{a,1\}=\{a,1\cdot1\}=2\{a,1\}\)，相减得 \(\{a,1\}=0\)，双线性再给出所有标量的括号为零。这个论证没有除以 \(2\)。

\ProseParagraphBreak
Poisson 理想正是能够取括号商的理想。若 \(I\) 满足定义中的条件，则
\[
	\{a+I,b+I\}=\{a,b\}+I
\]
不依赖代表：改动任一代表所增加的括号都在 \(I\) 中。三条恒等式由 \(A\) 下降到 \(A/I\)。反过来，若自然商同态保持某个商上的括号，\(b\in I\) 时 \(\{a,b\}\) 的像必须为零，所以 \(I\) 必为 Poisson 理想。特别地，与每个元素的括号均为零的元素组成\term[Poisson-zhongxin]{Poisson 中心}[Poisson center]，记为
\[
	Z_{\mathrm P}(A)=\{\,c\in A\mid \{c,a\}=0\ \text{对所有 }a\in A\,\}.
\]
乘积法则保证它是含幺子代数；若 \(c\in Z_{\mathrm P}(A)\)、\(\lambda\in k\)，则 \((c-\lambda)\) 是 Poisson 理想。
\symidx{\(Z_{\mathrm P}(A)\)：Poisson 代数的 Poisson 中心}

\begin{proposition}\label{b5:prop:standard-symplectic-poisson}
	设 \(k\) 为任意域，\(n\ge1\)，\(A=k[q_1,\ldots,q_n,p_1,\ldots,p_n]\)。公式
	\begin{equation}\label{b5:eq:standard-symplectic-poisson}
		\{f,g\}=\sum_{i=1}^n
		\left(\frac{\partial f}{\partial q_i}\frac{\partial g}{\partial p_i}
		-\frac{\partial f}{\partial p_i}\frac{\partial g}{\partial q_i}\right)
	\end{equation}
	使 \(A\) 成为 Poisson 代数。这是满足
	\[
		\{q_i,q_j\}=0,\qquad \{p_i,p_j\}=0,\qquad
		\{q_i,p_j\}=\delta_{ij}
	\]
	的唯一 Poisson 括号，称为\term[biaozhun-xin-kuohao]{标准辛括号}[standard symplectic Poisson bracket]。
\end{proposition}
\begin{proof}
	双线性及两变量的乘积法则由多项式求导得到。代入 \(g=f\) 后，每一对乘积之差都为零，故括号交替，包括特征二的情形。

	为核对 Jacobi，记三项和为 \(J(f,g,h)\)。只用两变量的乘积法则展开 \(J(f,g,uv)\)，得到
	\[
		J(f,g,uv)=J(f,g,u)v+uJ(f,g,v).
	\]
	确实，除右端外多出的项是
	\[
		\{g,u\}\bigl(\{f,v\}+\{v,f\}\bigr)
		+\bigl(\{f,u\}+\{u,f\}\bigr)\{g,v\}=0.
	\]
	由 \(J\) 的循环对称性，它在另两个变量上也满足乘积法则。因此只须检查三个变量都是代数生成元的情形：反复拆开单项式，就把一般的 \(J\) 写成这些值乘以多项式的和。生成元的两两括号是 \(0\) 或 \(\pm1\)，与任何元素的括号均为零，故所有这些值为零。于是 Jacobi 恒等式成立。

	若另一个括号具有所列生成元值，反复应用乘积法则，先把第一变量的单项式拆开，再把第二变量拆开，全部结果都由生成元的两两括号决定；双线性给出对全部多项式的唯一性。
\end{proof}

按 \((q_1,\ldots,q_n,p_1,\ldots,p_n)\) 排列，生成元括号的矩阵为
\[
	\begin{pmatrix}0&I_n\\-I_n&0\end{pmatrix}.
\]
它在任意特征下均可逆，且交替。在通常的全次数分次中，式\eqref{b5:eq:standard-symplectic-poisson}把次数分别为 \(r,s\ge1\) 的齐次多项式的括号送入次数 \(r+s-2\) 的齐次分量，允许结果为零；这与 PBW 括号下降一次不同。比如 \(n=1\) 时，非负整数 \(a,b,c,d\) 满足
\[
	\{q^ap^b,q^cp^d\}=(ad-bc)q^{a+c-1}p^{b+d-1}.
\]
若某个指数为 \(-1\)，对应系数必为零，该项按零解释。整数系数在 \(k\) 中取其像；例如 \(\{q^2p,qp^2\}=3q^2p^2\)，在特征三时为零。

\ProseParagraphBreak
固定 \(f\in A\)，记 \(H_f(g)=\{f,g\}\)。这是一个 \(k\) 导子，称为 \(f\) 的\term[Hamilton-daozi]{Hamilton 导子}[Hamiltonian derivation]。\footnote{哈密顿（William Rowan Hamilton，1805-1865），爱尔兰数学家、物理学家，研究四元数、几何光学及力学。}这里固定把 \(f\) 放在括号的第一位置。因此标准辛括号给出
\[
	H_f=\sum_i\left(
	\frac{\partial f}{\partial q_i}\frac{\partial}{\partial p_i}
	-\frac{\partial f}{\partial p_i}\frac{\partial}{\partial q_i}\right),
	\qquad H_{q_i}=\frac{\partial}{\partial p_i},\quad
	H_{p_i}=-\frac{\partial}{\partial q_i}.
\]
\symidx{\(H_f=\{f,-\}\)：采用第一变量约定的 Hamilton 导子}

\begin{proposition}\label{b5:prop:hamilton-derivation-commutator}
	设 \(A\) 为域 \(k\) 上的 Poisson 代数。赋值 \(f\mapsto H_f=\{f,-\}\) 是从 \(A\) 的 Lie 代数到 \(\operatorname{Der}_k(A)\) 的 Lie 同态；其核为 \(Z_{\mathrm P}(A)\)。具体地，对 \(f,g\in A\)，有
	\[
		[H_f,H_g]=H_{\{f,g\}}.
	\]
\end{proposition}
\begin{proof}
	线性及导子性质已由定义给出。对任意 \(a\in A\)，Jacobi 与反对称性给出
	\[
		[H_f,H_g](a)
		=\{f,\{g,a\}\}-\{g,\{f,a\}\}
		=\{\{f,g\},a\}.
	\]
	两端作用于每个元素都相等，即得同态性质；\(H_f=0\) 的定义正是 \(f\) 属于 Poisson 中心。
\end{proof}

例如在 \(k[q,p]\) 中，\(H_{qp}=p\partial/\partial p-q\partial/\partial q\)，于是
\[
	H_{qp}(q^ap^b)=(b-a)q^ap^b.
\]
即使某个元素不属于 Poisson 中心，也可能被某个固定的 Hamilton 导子零化。例如 \(qp\) 被 \(H_{qp}\) 零化，但 \(\{qp,q\}=-q\ne0\)。对 \(A=k[q_1,\ldots,q_n,p_1,\ldots,p_n]\) 的标准辛括号，检验与全部生成元的括号为零，就是检验所有偏导数为零，所以
\[
	Z_{\mathrm P}(A)=
	\begin{cases}
		k,&\operatorname{char}k=0,\\
		k[q_1^\ell,\ldots,q_n^\ell,p_1^\ell,\ldots,p_n^\ell],&\operatorname{char}k=\ell>0.
	\end{cases}
\]
公式可逐个单项式检查：一个变量的偏导为零，恰好要求每个非零项在该变量上的指数被特征整除；特征零时则要求该指数为零。此处没有对系数取根，所以不要求底域完美。

\ProseParagraphBreak
还可以直接比较这两种 Poisson 结构。在前面的 Heisenberg 例子中，\(z\) 属于 \(k[x,y,z]\) 的 Poisson 中心。取商 \(z=c\) 后，得到 \(k[x,y]\) 上的括号
\[
	\{f,g\}_c=c\left(
	\frac{\partial f}{\partial x}\frac{\partial g}{\partial y}
	-\frac{\partial f}{\partial y}\frac{\partial g}{\partial x}\right).
\]
取 \(c=0\) 时括号为零；取 \(c=1\) 并令 \(q=x,p=y\)，就是标准辛括号。这是把一个中央坐标赋值所得到的商，不是保持原全次数的分次同构：\(z-1\) 本来就不是齐次多项式。第四卷第13.4.4小节从形式星积的一次交换子得到 Poisson 括号，并在特征零下构造标准辛括号的星积；这里的括号本身和 Jacobi 验证适用于任意特征。

% * 新增小节状态：初稿已完成；原有内容保留。

% !TeX root = ../../Book5.tex
% 纲要标题：### 15.3 具体计算
% 纲要位置：工作记录/计划纲要.md，第 4564 行。

\section{具体计算}
\label{b5:sec:1503}


% 纲要范围：
%
% * 有序单项式
% * sl_2 的包络代数
% * Casimir 元的例子
% * 伴随局部幂零性传播生成向量的零化关系；Casimir 主符号的 Poisson 中心性与最高权标量对称性
%

有序单项式基使元素的坐标和乘积都可以直接计算。先固定次序，再按交换关系重排，并保留每次交换产生的低次项；漏掉括号项会改变元素。伴随导子也与这种重排相容，后面将用这一性质证明最高权商模上的局部幂零性。

\subsection{有序单项式}
\label{b5:subsec:1503:01}

设 \(\mathfrak g\) 有有序基 \(x_1<\cdots<x_n\)。PBW 使每个元素唯一写成有限和
\[
	u=\sum_{a_1,\ldots,a_n\ge0} c_{a_1,\ldots,a_n}x_1^{a_1}\cdots x_n^{a_n}.
\]
相邻逆序替换是求这些系数的算法；最高次数的部分按交换多项式相乘，而低次数由 Lie 括号补正。特别地，PBW 是向量空间基定理，不表示这些 \(x_i\) 在 \(U(\mathfrak g)\) 中已经交换。

\ProseParagraphBreak
若一个模由 \(v\) 生成，要把 \(x^bv=0\) 传到任意 \(uv\)，就须把作用在左端的 \(x\) 移过 \(u\)。记 \(D(u)=xu-ux\)，关系 \(xu=D(u)+ux\) 表明，每次移动或者产生一个伴随导子项，或者使 \(x\) 更接近生成向量。下面的命题用这两种零化关系控制整个展开。

\begin{proposition}\label{b5:prop:enveloping-adjoint-nilpotence}
	设 \(\mathfrak g\) 为域 \(k\) 上的 Lie 代数，\(x\in\mathfrak g\)，并假设 \(\operatorname{ad}x\) 在 \(\mathfrak g\) 上局部幂零，即对每个 \(y\in\mathfrak g\)，都存在整数 \(a\ge1\) 使 \((\operatorname{ad}x)^a y=0\)。则 \(\operatorname{ad}x\) 在泛包络代数 \(U(\mathfrak g)\) 上的导子延拓也局部幂零。在此假设下，若左 \(\mathfrak g\) 模 \(V\) 由向量 \(v\in V\) 生成，且 \(x^b v=0\) 对某个整数 \(b\ge1\) 成立，则 \(x\) 在整个 \(V\) 上局部幂零。
\end{proposition}
\begin{proof}
	记 \(D(u)=xu-ux\)。它是结合代数的导子，且在 \(\mathfrak g\) 上就是给定的伴随作用。对一个乘积 \(y_1\cdots y_r\)，选 \(a_i\ge1\) 使 \(D^{a_i}y_i=0\)。反复使用 Leibniz 规则，\(D^N\) 的每项是在各因子上分别作用 \(j_i\) 次、且 \(\sum_i j_i=N\) 的乘积。若 \(N>\sum_i(a_i-1)\)，至少一个 \(j_i\ge a_i\)，该项为零。线性组合只有有限项，故 \(D\) 在每个包络代数元素上局部幂零。 再由 \(xu=D(u)+ux\)，对 \(N\) 归纳得到
	\[
		x^Nu=\sum_{j=0}^N\binom NjD^j(u)x^{N-j}.
	\]
	若 \(D^a(u)=0\)、\(x^bv=0\)，取 \(N\ge a+b-1\)，每项或有 \(j\ge a\)，或有 \(N-j\ge b\)，于是 \(x^Nuv=0\)。\(V\) 中元素为有限个 \(uv\) 的线性组合，取各个界的最大值即可。
\end{proof}
这个证明不要求一个统一的幂同时零化整个包络代数；指数可以依赖具体元素。生成模也可以无限维，只需生成向量满足所列幂关系。

\ProseParagraphBreak
例如在 \(\mathfrak{sl}_2(\mathbb C)\) 中，取矩阵单位给出的标准基 \(e=E_{12}\)、\(f=E_{21}\)、\(h=E_{11}-E_{22}\)。对 \(D=\operatorname{ad}e\)，有
\[
	De=0,\qquad Dh=-2e,\qquad Df=h,\qquad D^2f=-2e,\qquad D^3f=0.
\]
因此 \(D^3\) 零化整个三维 Lie 代数，命题便保证 \(D\) 在其包络代数上局部幂零。如何利用 \(ev=0\) 把这一结论传到最高权生成模，将在 Verma 模定义以后具体说明。

\subsection{\texorpdfstring{\(\mathfrak{sl}_2\) 的包络代数}{sl2 的包络代数}}
\label{b5:subsec:1503:02}

对 \(\mathfrak{sl}_2(\mathbb C)\)，选择次序 \(f<h<e\)。则
\[
	\{\,f^a h^b e^c\mid a,b,c\ge0\,\}
\]
为 \(U(\mathfrak{sl}_2)\) 的基。三个基本重排为
\[
	ef=fe+h,\qquad eh=(h-2)e,\qquad hf=f(h-2).
\]
例如
\[
	efh=f(h-2)e+h^2,\qquad e^2f=fe^2+2(h-1)e.
\]
第二式由 \(e(ef)=e(fe+h)\) 展开，再把 \(eh\) 按规则重排得到。
\ProseParagraphBreak
\cref{b5:lem:sl2-ladder}的算子恒等式同样是包络代数中的恒等式：
\[
	ef^r=f^re+r f^{r-1}(h-r+1)\qquad(r\ge1).
\]
它的原证明只用了三个交换关系，故可以在当前结合代数里原样进行。以 \(r=2\) 继续计算，可得
\[
	\begin{aligned}
		e^2f^2 &=e\bigl(f^2e+2f(h-1)\bigr)\\
		&=f^2e^2+2f(h-1)e+2(fe+h)(h-1)\\
		&=f^2e^2+2f\bigl((h-1)+(h-3)\bigr)e+2h(h-1)\\
		&=f^2e^2+4f(h-2)e+2h(h-1).
	\end{aligned}
\]
从第二行到第三行用了 \(e(h-1)=(h-3)e\)，再合并两个含 \(fe\) 的项。末式已是所选次序下的表达。这里的低次项在具体表示上可能与高次项相消，但在包络代数本身并没有消失。

\subsection{Casimir 元的例子}
\label{b5:subsec:1503:03}

\begin{proposition}\label{b5:prop:sl2-casimir-central}
	记 \(\mathfrak{sl}_2=\mathfrak{sl}_2(\mathbb C)\)，\(E_{ij}\) 为矩阵单位，取标准基 \(e=E_{12}\)、\(f=E_{21}\)、\(h=E_{11}-E_{22}\)，并通过自然映射把它们看作泛包络代数的元素。在 \(U(\mathfrak{sl}_2)\) 中，元素
	\[
		C=h^2+2h+4fe=h^2-2h+4ef
	\]
	属于中心。对任意复 \(\mathfrak{sl}_2\) 模 \(V\)，若 \(v\in V\setminus\{0\}\)、\(\lambda\in\mathbb C\) 满足 \(ev=0\)、\(hv=\lambda v\)，则 \(C\) 在 \(v\) 上作用为 \(\lambda(\lambda+2)\)；如果整个模由 \(v\) 生成，同一个标量作用遍及整个模。
\end{proposition}
\symidx{\(C=h^2+2h+4fe\)：\(\mathfrak{sl}_2\) 包络代数中采用的 Casimir 元}
\begin{proof}
	两个表达式由 \(ef-fe=h\) 相等。\cref{b5:lem:sl2-casimir-operator}中的交换子计算只用这三条关系，因此在包络代数中仍给出 \([e,C]=[f,C]=[h,C]=0\)。这三个元素生成整个结合代数，所以 \(C\) 中央。最高向量 \(v\) 满足 \(ev=0\)、\(hv=\lambda v\)，代入第一式即得标量；对任意 \(u\)，中央性使 \(C(uv)=uCv=\lambda(\lambda+2)uv\)。
\end{proof}
PBW 还说明 \(C\) 不是一个标量元素：它的二次主符号为 \(q=h^2+4fe\)，在交换多项式环 \(\mathbb C[f,h,e]\) 中非零。\cref{b5:prop:pbw-poisson}把交换子下降一次所留下的部分变成 Poisson 括号；中央元素的主符号因而属于 Poisson 中心。在当前坐标中，这一点也可直接核对：
\[
	\begin{aligned}
		\{e,q\}&=-4he+4he=0,\\
		\{f,q\}&=4hf-4fh=0,\\
		\{h,q\}&=4\bigl((-2f)e+f(2e)\bigr)=0.
	\end{aligned}
\]
这里使用 \(\{e,f\}=h\)、\(\{h,e\}=2e\)、\(\{h,f\}=-2f\) 及交换乘法。三个生成元与 \(q\) 的括号均为零，再由 Leibniz 法则得到 \(q\) 与每个多项式的括号为零。\(C\) 在上述最高权生成模上按标量作用，并不意味着它在包络代数里本来就是标量。
\ProseParagraphBreak
若按 Killing 型取对偶基，则对应的 Casimir 元是 \(C/8\)，见\subsecref{b5:subsec:1403:04}末尾的归一化计算。选择哪一种写法并不改变中央性，但会改变其特征值常数。本编一般半单代数的公式将采用 Killing 归一化。

% !TeX root = ../../Book5.tex
% 纲要标题：### 15.4 诱导模初步
% 纲要位置：工作记录/计划纲要.md，第 4571 行。

\section{诱导模初步}
\label{b5:sec:1504}


% 纲要范围：
%
% * 子代数上的扩张标量
% * 三角分解的预告
% * 诱导模与 Verma 模
% * 比较 Cartan 与 Borel 诱导；最高权向量的 Hom 泛性质、左理想商、截断相容性与同一 Casimir 标量下的非分裂扩张
%
% ---
%

子代数的表示可以借扩张标量诱导为较大 Lie 代数的表示。与有限群的诱导不同，这样构造出的模通常无限维。PBW 给出所需的自由右模基，由此可构造最高权模，并证明其最高权向量非零。

\subsection{子代数上的扩张标量}
\label{b5:subsec:1504:01}

\begin{theorem}\label{b5:thm:lie-induction-adjunction}
设 \(\mathfrak h\subseteq\mathfrak g\) 为域 \(k\) 上 Lie 代数的子代数。自然映射
\(U(\mathfrak h)\to U(\mathfrak g)\) 单射，且 \(U(\mathfrak g)\) 是自由右 \(U(\mathfrak h)\) 模。对左 \(\mathfrak h\) 模 \(M\)，定义
\[
\operatorname{Ind}_{\mathfrak h}^{\mathfrak g}M
=U(\mathfrak g)\otimes_{U(\mathfrak h)}M,
\]
其中 \(\mathfrak g\) 通过第一因子左乘作用。对任意左 \(\mathfrak g\) 模 \(V\)，记 \(V|_{\mathfrak h}\) 为限制到 \(\mathfrak h\) 的模，则有自然同构
\[
\operatorname{Hom}_{\mathfrak g}
\left(\operatorname{Ind}_{\mathfrak h}^{\mathfrak g}M,V\right)
\cong
\operatorname{Hom}_{\mathfrak h}(M,V|_{\mathfrak h}).
\]
诱导函子保持正合列，且 \(M\to\operatorname{Ind}_{\mathfrak h}^{\mathfrak g}M\)、\(m\mapsto1\otimes m\) 单射。
\end{theorem}
\begin{proof}
把 \(\mathfrak h\) 的基扩为 \(\mathfrak g\) 的基，排序时令补充基元在前，子代数的基元在后。PBW 先说明子代数的有序单项式在大代数中仍独立，所以自然映射单射；又说明每个大代数基元唯一写成一个补充基元有序单项式乘以一个 \(U(\mathfrak h)\) 基元。这正给出自由右模基，其中包括空单项式一。
伴随同构把 \(F\) 送到 \(m\mapsto F(1\otimes m)\)。反向把 \(\mathfrak h\) 同态 \(T\) 送到 \(u\otimes m\mapsto u\cdot T(m)\)；\(T\) 与 \(U(\mathfrak h)\) 作用相容，保证公式对张量关系良定义。两者直接互逆并自然。
右自由性使诱导的底层向量空间为若干份 \(M\) 的直和，故保持正合性；空单项式对应的直和分量正好给出所列单射。
\end{proof}
例如从零子代数的平凡一维模诱导得到 \(U(\mathfrak g)\) 本身。只要 \(\mathfrak g\ne0\)，它就无限维。诱导的泛性质先构造所有可能的作用，有限维性往往要靠随后取商和加入额外关系获得。

\subsection{三角分解的动机}
\label{b5:subsec:1504:02}

\begin{proposition}\label{b5:prop:pbw-triangular-factorization}
设域 \(k\) 上的 Lie 代数 \(\mathfrak g\) 作为向量空间分解为
\[
\mathfrak g=\mathfrak n^-\oplus\mathfrak h\oplus\mathfrak n^+,
\]
其中三项都是子代数。则乘法给出向量空间同构
\[
U(\mathfrak n^-)\otimes_kU(\mathfrak h)\otimes_kU(\mathfrak n^+)
\longrightarrow U(\mathfrak g).
\]
若另有 \(\mathfrak b=\mathfrak h\oplus\mathfrak n^+\) 为子代数，则
\(U(\mathfrak g)\cong U(\mathfrak n^-)\otimes_kU(\mathfrak b)\)
作为右 \(U(\mathfrak b)\) 模相容。
\end{proposition}
\begin{proof}
分别取三项的有序基，并按负、零、正的次序合并。PBW 给出的有序字恰好是三段有序字的乘积，因此乘法把左边的张量积基双射到右边的基。若后两项合成子代数，按同一分组即可得到右模陈述。
\end{proof}
这个同构一般不是结合代数同构，因为三段之间的元素未必交换。它的作用是把自由方向和将要规定作用的方向分开。若 \(\mathfrak h\) 交换、\([\mathfrak h,\mathfrak n^+]\subseteq\mathfrak n^+\)，则 \(\mathfrak n^+\) 是 \(\mathfrak b\) 的理想，且由 \(\mathfrak h\) 交换及 \(\mathfrak n^+\) 为子代数得到
\[
[\mathfrak b,\mathfrak b]\subseteq\mathfrak n^+.
\]
将任意线性泛函 \(\lambda:\mathfrak h\to k\) 延伸到 \(\mathfrak b\)，使它在 \(\mathfrak n^+\) 上为零，就零化了 \([\mathfrak b,\mathfrak b]\)，因而给出 \(\mathfrak b\) 的一维模：\(\mathfrak h\) 按 \(\lambda\) 作用，\(\mathfrak n^+\) 按零作用。诱导以后，底层向量空间便是 \(U(\mathfrak n^-)\)。
\ProseParagraphBreak
在 \(\mathfrak{sl}_2\) 中可取
\(\mathfrak n^-=\mathbb Cf\)、\(\mathfrak h=\mathbb Ch\)、\(\mathfrak n^+=\mathbb Ce\)。
这正是次序 \(f<h<e\) 的来源。一般半单 Lie 代数中的三项，将在根分解中由负根、Cartan 子代数和正根构成；本命题只用给定的子空间分解，没有预先假定那套结构存在。

\subsection{诱导模与 Verma 模}
\label{b5:subsec:1504:03}

\begin{definition}\label{b5:def:sl2-verma}
设 \(\lambda\in\mathbb C\)，记 \(\mathfrak{sl}_2=\mathfrak{sl}_2(\mathbb C)\)，\(E_{ij}\) 为矩阵单位，\(e=E_{12}\)、\(f=E_{21}\)、\(h=E_{11}-E_{22}\)，令 \(\mathfrak b=\mathbb Ch\oplus\mathbb Ce\)。在一维复空间 \(\mathbb C_\lambda\) 上规定 \(h\) 作用为 \(\lambda\)、\(e\) 作用为零，得到一个左 \(\mathfrak b\) 模。诱导模
\[
M(\lambda)=U(\mathfrak{sl}_2)\otimes_{U(\mathfrak b)}\mathbb C_\lambda
\]
称为 \(\mathfrak{sl}_2\) 的最高权为 \(\lambda\) 的\term[Verma-mo]{Verma 模}[Verma module]。
\end{definition}
\symidx{\(M(\lambda)\)：最高权为 \(\lambda\) 的 Verma 模}
由 PBW 三角分解，\(M(\lambda)\) 有基
\(v_j=f^j\otimes1\)、\(j\ge0\)，并满足
\[
\begin{aligned}
hv_j&=(\lambda-2j)v_j,\qquad fv_j=v_{j+1}\qquad(j\ge0),\\
ev_0&=0,\qquad ev_j=j(\lambda-j+1)v_{j-1}\qquad(j\ge1).
\end{aligned}
\]
这些式子与有限模型相同，但这里没有预先规定 \(v_{m+1}=0\)。特别地，最高向量 \(v_0\) 非零；这由 PBW 证明，而不是诱导记号本身的形式性质。

\ProseParagraphBreak

若只保留 \(h\) 的一维作用，从 \(\mathbb Ch\) 诱导得到
\[
W(\lambda)=U(\mathfrak{sl}_2)\otimes_{U(\mathbb Ch)}\mathbb C_\lambda,
\]
这里把 \(\mathbb C_\lambda\) 限制为 \(\mathbb Ch\) 模。按次序 \(f<e<h\)，PBW 给出自由右 \(U(\mathbb Ch)\) 基 \(f^ae^b\)、\(a,b\ge0\)，故 \(W(\lambda)\) 有基 \(f^ae^b\otimes1\)。其生成向量 \(w=1\otimes1\) 满足 \(hw=\lambda w\)，但 \(ew=e\otimes1\ne0\)，所以只是给定权的向量。从 \(\mathfrak b\) 诱导时，还同时规定 \(e\) 零化生成向量，才得到只沿 \(f\) 方向自由生成的 Verma 模。

\begin{corollary}\label{b5:cor:sl2-verma-universal}
设 \(\lambda\in\mathbb C\)，\(E_{ij}\) 为二阶矩阵单位，\(e=E_{12}\)、\(f=E_{21}\)、\(h=E_{11}-E_{22}\) 为 \(\mathfrak{sl}_2(\mathbb C)\) 的标准基，\(M(\lambda)\) 为最高权 \(\lambda\) 的 Verma 模，记 \(v_0=1\otimes1\)、\(v_j=f^jv_0\)。对任意复左 \(\mathfrak{sl}_2\) 模 \(V\)，取生成向量的像给出自然的复向量空间同构
\[
\operatorname{Hom}_{\mathfrak{sl}_2}(M(\lambda),V)
\cong\{v\in V\mid ev=0,\ hv=\lambda v\}.
\]
右侧包含零向量。对应于 \(v\) 的同态把每个 \(v_j\) 送到 \(f^jv\)，其像为 \(v\) 生成的子模；特别地，若 \(v\) 生成 \(V\)，则该同态满射。
\end{corollary}
\begin{proof}
由\cref{b5:thm:lie-induction-adjunction}，左侧自然同构于 \(\operatorname{Hom}_{\mathfrak b}(\mathbb C_\lambda,V|_{\mathfrak b})\)。一个这样的同态由 \(1\) 的像 \(v\) 唯一确定，与 \(e,h\) 的作用相容恰好等于 \(ev=0\)、\(hv=\lambda v\)，不要求 \(v\ne0\)。伴随同构的反向公式是 \(u\otimes1\mapsto uv\)，故 \(v_j\mapsto f^jv\)。其像就是 \(U(\mathfrak{sl}_2)v\)，这证明关于像和满射的陈述。
\end{proof}

由\cref{b5:prop:enveloping-adjoint-nilpotence}及 \(\operatorname{ad}e\) 在 \(\mathfrak{sl}_2\) 上的幂零性，\(ev_0=0\) 已足以推出 \(e\) 在 \(M(\lambda)\) 上局部幂零。基上的作用公式也直接给出 \(e^{j+1}v_j=0\)；有限线性组合取这些界的最大值即可。另一方面，\(f^rv_0=v_r\ne0\) 对每个 \(r\ge0\) 成立，所以 \(f\) 在 Verma 模上并不局部幂零。

\ProseParagraphBreak
诱导时的一维关系只规定生成元 \(v_0\) 满足 \(ev_0=0\)、\(hv_0=\lambda v_0\)，并没有规定 \(e\) 在整个诱导模上为零。例如
\[
e(fv_0)=[e,f]v_0+f(ev_0)=\lambda v_0.
\]
也可把此模写成左模商
\[
M(\lambda)\cong
U(\mathfrak{sl}_2)\big/
\bigl(U(\mathfrak{sl}_2)e+U(\mathfrak{sl}_2)(h-\lambda)\bigr).
\]
分母是生成元关系产生的左理想，商中只保留左作用，不是结合代数的商。具体地，映射 \(u\mapsto uv_0\) 零化这一左理想；按 \(f<h<e\) 的 PBW 次序，含正次 \(e\) 的项被消去，而 \(f^a h^b\) 模该理想等于 \(\lambda^b f^a\)。因此商由 \(f^a\) 的类生成，它们的像是已经证明独立的 \(v_a\)，这同时证明上述同构。

\ProseParagraphBreak

若把 \(e\)、\(h-\lambda\) 改为生成双边理想 \(J_\lambda\)，商代数就会塌缩。在该商中 \(e=0\)，所以 \(h=[e,f]=0\)，同时 \(h=\lambda\cdot1\)。当 \(\lambda\ne0\) 时，单位为零，商是零代数。当 \(\lambda=0\) 时，再由 \([h,f]=-2f\) 得 \(f=0\)，故商只剩标量。将 \(e,f,h\) 都送到零的增广同态说明单位没有消去，所以此时商恰为 \(\mathbb C\)。双边关系会限制整个代数，左理想关系则只规定一个生成向量的作用。

\ProseParagraphBreak

取 \(\lambda=2\)，前几个权及升降系数可直接画成
\[
\begin{tikzcd}[column sep=2.4em]
\substack{v_0\\2}\arrow[r,shift left=1.2ex,"f"]
&\substack{v_1\\0}\arrow[l,shift left=1.2ex,"e:2"]\arrow[r,shift left=1.2ex,"f"]
&\substack{v_2\\-2}\arrow[l,shift left=1.2ex,"e:2"]\arrow[r,shift left=1.2ex,"f"]
&\substack{v_3\\-4}\arrow[l,shift left=1.2ex,"e:0"]\arrow[r,shift left=1.2ex,"f"]
&\substack{v_4\\-6}\arrow[l,shift left=1.2ex,"e:-4"]\arrow[r,"f"]
&\cdots
\end{tikzcd}
\]
每支 \(f\) 箭的系数都是 \(1\)，\(e:c\) 表示升算子的系数为 \(c\)。尤其 \(fv_2=v_3\ne0\)，但 \(ev_3=0\)。所以前三个权的线性生成空间不是子模；尾部 \(\operatorname{span}\{v_3,v_4,\ldots\}\) 才是子模。将尾部取商，才使从权 \(-2\) 继续下降的箭变为零，得到有限维的 \(L_2\)。下面证明这一尾部是唯一的非零真子模。

\ProseParagraphBreak
有限截断须与升算子相容。在 \(M(\lambda)\) 的一个非零商模中，记生成向量的像为 \(\bar v_0\ne0\)。若下降链终止，取使 \(f^N\bar v_0=0\) 成立的最小正整数 \(N\)，则 \(f^{N-1}\bar v_0\ne0\)。作用 \(e\) 得
\[
0=ef^N\bar v_0=N(\lambda-N+1)f^{N-1}\bar v_0,
\]
因此 \(\lambda=N-1\)。这里使用了最小零化指数；更大的幂也可能零化生成向量，却不能代入这一推论。例如 \(\lambda=\dfrac12\) 时强加 \(f^2\bar v_0=0\)，便由 \(ef^2\bar v_0=-f\bar v_0\) 得 \(f\bar v_0=0\)，再由 \(ef\bar v_0=\dfrac12\bar v_0\) 得 \(\bar v_0=0\)，只会得到零模。

\begin{theorem}\label{b5:thm:sl2-verma-submodules}
设 \(\lambda\in\mathbb C\)，\(M(\lambda)\) 为复 \(\mathfrak{sl}_2(\mathbb C)\) 的最高权为 \(\lambda\) 的 Verma 模，\(v_\lambda\) 为其标准最高权生成元，记 \(v_j=f^jv_\lambda\)，其中 \(f=E_{21}\)、\(E_{21}\) 为矩阵单位。若 \(\lambda\notin\mathbb Z_{\ge0}\)，则 \(M(\lambda)\) 简单。若 \(\lambda=m\in\mathbb Z_{\ge0}\)，则它唯一的非零真子模为
\[
N_m=\operatorname{span}_{\mathbb C}\{v_{m+1},v_{m+2},\ldots\}
\cong M(-m-2),
\]
且 \(M(m)/N_m\cong L_m\)，其中 \(L_m\) 为最高权 \(m\)、维数 \(m+1\) 的不可约复 \(\mathfrak{sl}_2\) 模。列
\(0\to N_m\to M(m)\to L_m\to0\) 不分裂。
\end{theorem}
\begin{proof}
任一非零子模中的非零向量只有有限多个权分量，而 \(h\) 的权两两不同。若该向量的 \(v_j\) 系数非零，就取多项式 \(p\)，使它在权 \(\lambda-2j\) 处取一，在其余出现的权处取零。子模被 \(h\) 的多项式保持，故 \(p(h)\) 将该向量投影为非零的 \(v_j\) 倍数；除去这个系数，便使子模包含 \(v_j\)。如果 \(\lambda\) 不是非负整数，所有升算子系数 \(r(\lambda-r+1)\)、\(r\ge1\) 都非零，所以能升到 \(v_0\)，再用 \(f\) 生成全部模，证明简单。
若 \(\lambda=m\)，则 \(ev_{m+1}=0\)、\(hv_{m+1}=(-m-2)v_{m+1}\)，所列尾部空间稳定。由\cref{b5:cor:sl2-verma-universal}，存在模同态 \(M(-m-2)\to N_m\)，把标准最高权生成元 \(w_0\) 送到 \(v_{m+1}\)。记源模的 PBW 基为 \(w_r=f^rw_0\)、\(r\ge0\)，这个同态具体为
\[
M(-m-2)\longrightarrow N_m,\qquad
w_r\longmapsto v_{m+1+r}
\]
它把源模的全部基向量双射到尾部的全部基向量，因此同时单射且满射，得到所述模同构。由第一种情形，\(M(-m-2)\) 简单。任何真子模不能包含 \(v_j\)、\(j\le m\)，否则可升到 \(v_0\)；若它非零，就包含某个 \(j\ge m+1\) 的 \(v_j\)，再升到 \(v_{m+1}\)，因而包含整个 \(N_m\)。而商的前 \(m+1\) 个基向量给出 \(L_m\)，故没有其他非零真子模。
最后，\(f\) 在 \(M(m)\) 上单射，在 \(L_m\) 上满足 \(f^{m+1}=0\)。若存在模同态截面，其像 \(S\) 是与 \(L_m\) 同构的非零有限维子模。由于 \(f\) 在整个 \(M(m)\) 上单射，限制 \(f|_S\) 也单射；由于截面与作用相容，又有 \((f|_S)^{m+1}=0\)。这与非零空间上的幂零算子不可能单射矛盾。因此列不分裂。
\end{proof}
在有限维商 \(L_m\) 上，升降公式给出 \(e^{m+1}=f^{m+1}=0\)，两个算子都幂零。这个结论包含了截断以后新得到的 \(f\) 的零化关系；无限维 Verma 模本身只有前面证明的 \(e\) 的局部幂零性。

\ProseParagraphBreak
Casimir 标量也不能把上述扩张分开。记 \(c(\lambda)=\lambda(\lambda+2)\)，则
\[
c(-\lambda-2)=(-\lambda-2)(-\lambda)=c(\lambda).
\]
由\cref{b5:prop:sl2-casimir-central}，\(C\) 在整个 \(M(m)\) 上就是标量 \(c(m)\)，在其子模 \(M(-m-2)\) 和商模 \(L_m\) 上也是同一个标量。例如 \(m=2\) 时，三者都取值 \(8\)。因此任何 \(C\) 的多项式也只能在三者上取同一个标量，无法给出区分两项的特征空间投影；而该列已经独立证明不分裂。在有限维简单模的非负整数最高权中，\(c(m)\) 两两不同；一旦允许上述无限维模，最高权 \(m\) 与 \(-m-2\) 便有相同的 Casimir 值。这说明 Casimir 不能分离这个扩张。

\ProseParagraphBreak
这个例子也说明有限维条件的重要性：有限维复 \(\mathfrak{sl}_2\) 模完全可约，并不意味着所有包络代数模都完全可约。后面的一般最高权构造会采用相同方法，先建立非零的诱导模，再证明何时其简单商有限维。


\begin{chaptersummary}
包络代数的定义只加入 \(xy-yx=[x,y]\) 所要求的关系。泛性质把 Lie 代数作用唯一延拓为结合代数作用，但不能单独保证原生成元仍然独立。PBW 定理通过重排的终止性、Jacobi 控制的相容性及正规式的独立性，才完成这一点。

\ProseParagraphBreak
标准滤过使包络代数的最高次乘法成为交换乘法；伴随分次与对称代数同构，交换子下降一次所产生的 Poisson 括号仍记录原 Lie 括号。主符号可用于增长与零因子判断，对称化则在特征零时给出与伴随作用相容的向量空间同构，乘法中的低次修正一般仍然存在。增广的一阶商 \(I/I^2\) 测量 Lie 代数的交换化，与 PBW 的 \(F_1/F_0\) 有不同含义。

\ProseParagraphBreak
子代数的包络代数嵌入较大的包络代数，PBW 进一步给出自由右模基，所以扩张标量是正合的。从 \(\mathfrak{sl}_2\) 的 Verma 模出发的同态，恰由目标模中具有给定权且被升算子零化的向量决定。Verma 模由无限下降链构成；只有在非负整数最高权处，才能相容地截断而得到有限维简单商。\chapref{b5:ch:18}将把这种构造推广到一般复半单 Lie 代数。

\ProseParagraphBreak
群代数 \(kG\) 中的群元素 \(g\) 在张量积两因子上同时作用，包络代数 \(U(\mathfrak g)\) 中的 \(x\in\mathfrak g\) 则按 \(x(v\otimes w)=xv\otimes w+v\otimes xw\) 作用。\chapref{b5:ch:21}将用余乘法 \(\Delta(g)=g\otimes g\) 和 \(\Delta(x)=x\otimes1+1\otimes x\) 同时描述这两条规则。
\end{chaptersummary}

\BookExercises

\begin{exercise}\label{b5:ex:15-affine-enveloping}
设 Lie 代数 \(\mathfrak g\) 有基 \(e,h\)，满足 \([h,e]=e\)，底域任意。
写出 \(U(\mathfrak g)\) 的生成关系与 PBW 基，并证明
\(he^m=e^m(h+m)\)。
\end{exercise}
\begin{solution}
生成关系为 \(he-eh=e\)。取 \(e<h\)，PBW 基为 \(e^ah^b\)、\(a,b\ge0\)。
由 \(he=e(h+1)\)，对 \(m\) 归纳：
\(he^{m+1}=e^m(h+m)e=e^{m+1}(h+m+1)\)。
整数 \(m\) 在底域中取其像，所以公式也适用于正特征。
\end{solution}

\begin{exercise}\label{b5:ex:15-heisenberg-order}
设 \(\mathfrak h=kx\oplus ky\oplus kz\) 为域 \(k\) 上满足 \([x,y]=z\)、\([x,z]=[y,z]=0\) 的 Heisenberg Lie 代数。在 \(x<y<z\) 的 PBW 次序中求 \(y^2x^2\)。
\end{exercise}
\begin{solution}
先由 \(yx=xy-z\) 得 \(yx^2=x^2y-2xz\)。
再左乘 \(y\)，先写成 \(y^2x^2=(yx^2)y-2(yx)z\)，再分别代入 \(yx^2=x^2y-2xz\) 与 \(yx=xy-z\)，利用 \(z\) 中央：
\[
y^2x^2=(x^2y-2xz)y-2(xy-z)z
=x^2y^2-4xyz+2z^2.
\]
该式保留了重排产生的一次、二次括号贡献。
\end{solution}

\begin{exercise}\label{b5:ex:15-one-dimensional-characters}
设 \(\mathfrak g=ke\oplus kh\) 为域 \(k\) 上满足 \([h,e]=e\) 的 Lie 代数。求它的全部一维表示。
再在二维空间上规定
\[
\rho(e)=\begin{pmatrix}0&b\\0&0\end{pmatrix},\qquad
\rho(h)=\begin{pmatrix}a&0\\0&c\end{pmatrix}.
\]
求 \(a,b,c\in k\) 必须满足的条件，并判断由第一坐标轴给出的短正合列何时分裂。
\end{exercise}
\begin{solution}
由 \([\mathfrak g,\mathfrak g]=ke\)，一维表示恰为 \(e\) 作用为零、\(h\) 作用为任意标量 \(a\) 的表示，记为 \(k_a\)。
矩阵计算给出 \([\rho(h),\rho(e)]=(a-c)\rho(e)\)，所以所需条件为 \((a-c-1)b=0\)。
第一坐标轴与商分别是 \(k_a\)、\(k_c\)，得到 \(0\to k_a\to V\to k_c\to0\)。若 \(b=0\)，第二坐标轴给出截面。
若 \(b\ne0\)，则 \(a=c+1\)，且 \(e\) 在 \(V\) 上非零；若列分裂，\(V\) 应是两个一维模的直和，而 \(e\) 在两项上都为零，矛盾。因此列恰在 \(b=0\) 时分裂。
\end{solution}

\begin{exercise}\label{b5:ex:15-enveloping-abelianization}
令 \(J\) 为 \(U(\mathfrak g)\) 中由全部交换子 \(uv-vu\) 生成的双边理想。
证明
\[
U(\mathfrak g)/J\cong\operatorname{Sym}\bigl(\mathfrak g/[\mathfrak g,\mathfrak g]\bigr).
\]
\end{exercise}
\begin{solution}
商代数交换，其中 \([x,y]=xy-yx=0\)，故来自
\(\mathfrak g\) 的映射通过其交换化。对称代数的泛性质给出右侧到左侧的同态。
反过来，\(\mathfrak g\to\mathfrak g/[\mathfrak g,\mathfrak g]\)
再到其对称代数的映射保持 Lie 括号，延拓到 \(U(\mathfrak g)\) 后零化所有结合交换子。
两复合在生成元及单位上为恒等，所以互逆。
\end{solution}

\begin{exercise}\label{b5:ex:15-direct-sum-enveloping}
证明
\(U(\mathfrak g\oplus\mathfrak l)\cong U(\mathfrak g)\otimes_kU(\mathfrak l)\)。
说明为什么直和中两因子的括号为零恰是所需条件。
\end{exercise}
\begin{solution}
把 \(x\in\mathfrak g\) 送到 \(x\otimes1\)，把 \(y\in\mathfrak l\) 送到
\(1\otimes y\)，两类像交换，故定义 Lie 代数同态并延拓为包络代数同态。
反向，两份包络代数在 \(U(\mathfrak g\oplus\mathfrak l)\) 中的像互相交换，
所以乘法给出张量积到它的同态。两复合在两份生成元及单位上为恒等。
若两因子的括号不为零，张量积中的像却总交换，就会丢掉该括号信息。
\end{solution}

\begin{exercise}\label{b5:ex:15-faithful-lie-not-enveloping}
设非零有限维 Lie 代数 \(\mathfrak g\) 有一个非零有限维表示 \(V\)。
证明延拓 \(U(\mathfrak g)\to\operatorname{End}_k(V)\) 不可能单射。
用 \(\mathfrak{sl}_2\) 自然模说明这不妨碍 Lie 代数作用忠实。
\end{exercise}
\begin{solution}
PBW 使 \(U(\mathfrak g)\) 无限维，而目标有限维，所以线性映射有非零核。
自然模中的 \(\mathfrak{sl}_2\) 矩阵表示显然忠实，但 \(e^2\) 在该模上作用为零；
PBW 又说明 \(e^2\ne0\) 于包络代数。因此两个忠实性问题不同。
\end{solution}

\begin{exercise}\label{b5:ex:15-augmentation-cotangent}
设 \(\mathfrak h=kx\oplus ky\oplus kz\) 为域 \(k\) 上满足 \([x,y]=z\)、\([x,z]=[y,z]=0\) 的 Heisenberg Lie 代数，\(I\) 为 \(U(\mathfrak h)\) 的增广理想。
令 \(N\subseteq M_3(k)\) 为严格上三角矩阵空间。求全部含幺代数同态
\(\Phi:U(\mathfrak h)\to kI_3+N\)，要求再模掉 \(N\) 后等于增广映射。
说明 \(z\in I^2\) 为何仍可能有非零的 \(\Phi(z)\)，并计算将目标进一步模掉 \(N^2\) 后保留的参数。
\end{exercise}
\begin{solution}
设 \(E_{ij}\) 为矩阵单位。任意选取六个标量，写成
\[
X=aE_{12}+bE_{23}+cE_{13},\qquad
Y=dE_{12}+sE_{23}+tE_{13}.
\]
有 \([X,Y]=(as-bd)E_{13}\)。又因 \(E_{13}N=NE_{13}=0\)，此交换子与 \(X,Y\) 都交换。
所以 \(x\mapsto X\)、\(y\mapsto Y\)、\(z\mapsto(as-bd)E_{13}\) 保持全部 Lie 括号，由泛性质唯一延拓；反过来，任一所求同态都必须这样给出。
例如取 \(X=E_{12}\)、\(Y=E_{23}\)，就有 \(\Phi(z)=E_{13}\ne0\)。\(z\in I^2\) 只保证它的像属于 \(N^2=kE_{13}\)，并不保证像为零。
模掉 \(N^2\) 后，\(c,t\) 及 \(z\) 的像消失，只留下 \(a,b,d,s\) 四个参数。这正是从二维空间 \(I/I^2\cong k\bar x\oplus k\bar y\) 到二维平方零空间 \(N/N^2\) 的线性映射。
\end{solution}

\begin{exercise}\label{b5:ex:15-perfect-augmentation}
设 \(\mathfrak g=\mathfrak{sl}_2(\mathbb C)\oplus\mathbb Ct\)，其中 \(t\) 中央。记 \(I\subseteq U(\mathfrak g)\) 为增广理想，\(K\) 为由 \(\mathfrak{sl}_2\) 生成的双边理想。
求 \(I/I^2\)、各 \(I^r\)、\(r\ge1\)，以及 \(\bigcap_{r\ge1}I^r\)。据此判断 \(I/I^2\ne0\) 是否保证增广理想滤过分离，并与 PBW 的次数一部分比较。
\end{exercise}
\begin{solution}
\(\mathfrak g/[\mathfrak g,\mathfrak g]\cong\mathbb Ct\)，由\cref{b5:prop:enveloping-augmentation-cotangent}，有 \(I/I^2\cong\mathbb Ct\)。
由 \(h=[e,f]\)、\(e=\frac12[h,e]\)、\(f=-\frac12[h,f]\)，\(\mathfrak{sl}_2\) 的每个元素都是这一子代数内括号的线性组合。由此归纳证明 \(K\subseteq I^r\)：\(r=1\) 显然；若 \(e,f,h\in I^r\)，则上述括号在 \(I^{2r}\subseteq I^{r+1}\) 中，而 \(I^{r+1}\) 是双边理想，故也包含它们生成的 \(K\)。
再由\cref{b5:prop:enveloping-quotient}，商代数为 \(U(\mathfrak g/\mathfrak{sl}_2)\cong U(\mathbb Ct)=\mathbb C[t]\)。商映射 \(\pi:U(\mathfrak g)\to U(\mathfrak g)/K\cong\mathbb C[t]\) 把 \(I\) 送到 \((t)\)，所以
\[
I^r=\pi^{-1}\bigl((t^r)\bigr)=K+t^rU(\mathfrak g).
\]
这里反向包含使用了 \(K\subseteq I^r\) 及 \(t^r\in I^r\)。一个非零多项式不能被任意高次 \(t^r\) 整除，因此 \(\bigcap_rI^r=K\ne0\)，滤过不分离。
PBW 的 \(F_1/F_0\cong\mathfrak g\) 是四维空间，增广的一阶商则只剩一维；非零的一阶商也不能使已经落入所有 \(I^r\) 的三个 Lie 方向重新显现。
\end{solution}

\begin{exercise}\label{b5:ex:15-filtered-poisson}
在域 \(k\) 上的多项式代数 \(k[x,y,z]\) 中，取 Heisenberg Lie 代数诱导的 Poisson 括号
\(\{x,y\}=z\)、\(\{x,z\}=\{y,z\}=0\)。计算 \(\{xy,x+y\}\)。
求 Poisson 中心，即与每个多项式的括号都为零的元素所成子代数；分别讨论特征零和特征 \(p>0\)。
\end{exercise}
\begin{solution}
Leibniz 法则给出 \(\{xy,x+y\}=-xz+yz=z(y-x)\)。对任意 \(P\in k[x,y,z]\)，有
\[
\{P,x\}=-z\frac{\partial P}{\partial y},\qquad
\{P,y\}=z\frac{\partial P}{\partial x},\qquad \{P,z\}=0.
\]
与三个生成元的括号为零，就由 Leibniz 法则推出与每个多项式的括号为零。多项式环中 \(z\) 不是零因子，所以条件恰是关于 \(x,y\) 的两个偏导数均为零。
特征零时，逐个单项式求导可知 \(P\) 只能含 \(z\)，中心是 \(k[z]\)。特征 \(p\) 时，关于 \(x\) 的导数为零恰要求每个非零项的 \(x\) 指数被 \(p\) 整除，\(y\) 同理，所以中心是 \(k[x^p,y^p,z]\)。这一结论不要求 \(k\) 完美，因为没有对系数取 \(p\) 次根。
\end{solution}

\begin{exercise}\label{b5:ex:15-pbw-hilbert}
求 \(\dim F_dU(\mathfrak{sl}_2)\) 及
\(\dim F_d/F_{d-1}\)，并写出伴随分次的 Hilbert 级数。
\end{exercise}
\begin{solution}
有序单项式为 \(f^ah^be^c\)，所以
\[
\dim F_d=\binom{d+3}{3},\qquad
\dim F_d/F_{d-1}=\binom{d+2}{2}.
\]
三个指数独立给出生成函数
\(\sum_{d\ge0}\binom{d+2}{2}t^d=(1-t)^{-3}\)。
\end{solution}

\begin{exercise}\label{b5:ex:15-symmetrization-cubic}
设底域特征零。将 \(\operatorname{sym}(x^2y)\) 写成
\(x^2y\) 与含 Lie 括号的较低次项之和。
\end{exercise}
\begin{solution}
六次排列中每个不同词出现两次，所以
\(\operatorname{sym}(x^2y)=(x^2y+xyx+yx^2)/3\)。
令 \(z=[x,y]\)，则 \(xyx=x^2y-xz\)，
\(yx^2=x^2y-xz-zx\)。
因此
\[
\operatorname{sym}(x^2y)
=x^2y-x[x,y]+\frac13\bigl[x,[x,y]\bigr].
\]
最高次仍是交换单项式 \(x^2y\)，其余项分别降低到次数二和次数一。
\end{solution}

\begin{exercise}\label{b5:ex:15-positive-characteristic-average}
在特征二的一维交换 Lie 代数中，说明未除以阶乘的二次“对称化和”
\(x^2\mapsto xx+xx\) 为什么不能代替正文的对称化同构。
\end{exercise}
\begin{solution}
该和等于 \(2x^2=0\)，而 \(x^2\) 在对称代数与包络代数中都非零。
所以未归一化的平均在次数二已经不是单射。
这里 PBW 仍成立，失败的是依赖除以 \(2!\) 的特定对称化公式。
\end{solution}

\begin{exercise}\label{b5:ex:15-sl2-order-e2f3}
在 \(\mathfrak{sl}_2\) 的 PBW 次序 \(f<h<e\) 中求 \(e^2f^3\)。
\end{exercise}
\begin{solution}
先用 \(ef^3=f^3e+3f^2(h-2)\)，再用 \(ef^2=f^2e+2f(h-1)\)：
\[
\begin{aligned}
e^2f^3
&=f^3e^2+3f^2(h-2)e
+3\bigl(f^2e+2f(h-1)\bigr)(h-2)\\
&=f^3e^2+3f^2\bigl((h-2)+(h-4)\bigr)e+6f(h-1)(h-2)\\
&=f^3e^2+6f^2(h-3)e+6f(h-1)(h-2).
\end{aligned}
\]
其中 \(e(h-2)=(h-4)e\)，不能在这一步把 \(e\) 与 \(h\) 当作交换。
\end{solution}

\begin{exercise}\label{b5:ex:15-casimir-transcendence}
在 \(U\bigl(\mathfrak{sl}_2(\mathbb C)\bigr)\) 中令 \(C=h^2+2h+4fe\)。
证明任意非零多项式 \(p(t)\in\mathbb C[t]\) 都满足 \(p(C)\ne0\)。
\end{exercise}
\begin{solution}
若 \(p\) 的次数为 \(r\)、首项系数为 \(a\ne0\)，则 \(p(C)\) 的最高次主符号为
\(a(h^2+4fe)^r\)，次数 \(2r\)。
PBW 将它置于多项式整环 \(\mathbb C[f,h,e]\)，该元素非零。
因此 \(p(C)\) 不可能为零，得到一个多项式代数 \(\mathbb C[C]\) 嵌入中心。
\end{solution}

\begin{exercise}\label{b5:ex:15-enveloping-antiautomorphism}
设 \(\mathfrak g\) 为域 \(k\) 上的 Lie 代数，\(V\) 为左 \(U(\mathfrak g)\) 模，\(S\) 为包络代数的反自同构，满足 \(S(x)=-x\) 对每个 \(x\in\mathfrak g\) 成立。证明
\[
(u\cdot\varphi)(v)=\varphi\bigl(S(u)v\bigr)
\quad\bigl(u\in U(\mathfrak g),\ \varphi\in V^*,\ v\in V\bigr)
\]
给出 \(V^*\) 上的左模作用，限制到 \(\mathfrak g\) 后正是对偶表示。
对 \(\mathfrak g=\mathfrak{sl}_2(\mathbb C)\)、\(V=\mathbb C^2\)，用 \(B(v,w)=v_1w_2-v_2w_1\) 构造一个 \(U(\mathfrak g)\) 模同构 \(V\to V^*\)。
\end{exercise}
\begin{solution}
由 \(S(ua)=S(a)S(u)\)，有
\[
\bigl(u\cdot(a\cdot\varphi)\bigr)(v)
=\varphi\bigl(S(a)S(u)v\bigr)
=\bigl((ua)\cdot\varphi\bigr)(v).
\]
单位作用为恒等，故是左模。\(x\in\mathfrak g\) 时，公式变成 \((x\cdot\varphi)(v)=-\varphi(xv)\)，即对偶表示。
若 \(x\) 为迹零的二阶矩阵，直接展开得 \(B(xv,w)+B(v,xw)=0\)。于是 \(T(v)=B(v,-)\) 满足
\(T(xv)(w)=-B(v,xw)=(x\cdot T(v))(w)\)，是 Lie 代数模同态，也就与整个包络代数作用相容。
\(B\) 的矩阵为 \(\left(\begin{smallmatrix}0&1\\-1&0\end{smallmatrix}\right)\)，可逆，故 \(T\) 为所求模同构。
\end{solution}

\begin{exercise}\label{b5:ex:15-enveloping-diagonal}
设 \(\mathfrak g=\mathfrak{sl}_2(\mathbb C)\)，\(e,f,h\) 为其标准基，\(C=h^2+2h+4fe\)，\(\Delta\) 为包络代数的余乘法，满足 \(\Delta(x)=x\otimes1+1\otimes x\) 对每个 \(x\in\mathfrak g\) 成立。
计算 \(\Delta(C)\)。再令 \(V=\mathbb C^2\) 为自然模，\(\tau(v\otimes w)=w\otimes v\)，证明 \(C\) 在 \(V\otimes V\) 上的作用为 \(4(1+\tau)\)，并写出其两个特征空间的投影。
\end{exercise}
\begin{solution}
利用 \(\Delta\) 保持乘法，展开得
\[
\Delta(C)=C\otimes1+1\otimes C
+2h\otimes h+4(f\otimes e+e\otimes f).
\]
自然模是 \(L_1\)，由\cref{b5:lem:sl2-casimir-operator}，\(C\) 为标量 \(1(1+2)=3\)。设标准基为 \(u_+,u_-\)，满足
\(hu_\pm=\pm u_\pm\)、\(eu_-=u_+\)、\(fu_+=u_-\)、\(eu_+=fu_-=0\)。记
\(A=h\otimes h+2(f\otimes e+e\otimes f)\)，则
\[
\begin{aligned}
A(u_\pm\otimes u_\pm)&=u_\pm\otimes u_\pm,\\
A(u_+\otimes u_-)&=-u_+\otimes u_-+2u_-\otimes u_+,\\
A(u_-\otimes u_+)&=-u_-\otimes u_++2u_+\otimes u_-.
\end{aligned}
\]
这四个基向量上的等式给出 \(A=2\tau-1\)。因此 \(\Delta(C)\) 的作用为 \(6\cdot1+2(2\tau-1)=4(1+\tau)\)。
对称张量与交错张量的特征值分别为 \(8\)、\(0\)，投影是
\[
P_8=\frac{1+\tau}{2}=\frac{\rho_{V\otimes V}(C)}8,\qquad
P_0=\frac{1-\tau}{2}=1-P_8.
\]
两个因子上的 Casimir 标量之和只给出 \(6\)；张量作用还包含余乘法产生的交叉项，它们使两个直和项具有不同的特征值。
\end{solution}

\begin{exercise}\label{b5:ex:15-strict-induction-infinite}
设 \(\mathfrak h\subsetneq\mathfrak g\) 为域上 Lie 代数，\(M\ne0\) 为 \(\mathfrak h\) 模。
证明 \(\operatorname{Ind}_{\mathfrak h}^{\mathfrak g}M\) 必无限维，即使 \(M\) 一维也如此。
\end{exercise}
\begin{solution}
把子代数基扩成大代数基，补充基元至少有一个 \(x\)。
PBW 右自由模基包含互不相同的 \(1,x,x^2,\ldots\)。
张量 \(M\) 后，它们给出无穷多份非零的 \(M\) 直和，
所以诱导模无限维。这个结论与底域特征无关。
\end{solution}

\begin{exercise}\label{b5:ex:15-induction-transitivity}
设 \(\mathfrak a\subseteq\mathfrak h\subseteq\mathfrak g\) 为子代数，
\(M\) 为 \(\mathfrak a\) 模。构造并验证诱导的传递同构。
\end{exercise}
\begin{solution}
同构为
\[
U(\mathfrak g)\otimes_{U(\mathfrak h)}
\bigl(U(\mathfrak h)\otimes_{U(\mathfrak a)}M\bigr)
\longrightarrow U(\mathfrak g)\otimes_{U(\mathfrak a)}M,
\quad u\otimes(v\otimes m)\longmapsto uv\otimes m.
\]
两层平衡关系均由结合律保持，所以良定。
逆为 \(u\otimes m\mapsto u\otimes(1\otimes m)\)，
其平衡性来自 \(a\otimes m=1\otimes am\) 于内层张量积。
两个复合利用外层平衡关系化为恒等，且都保持左 \(\mathfrak g\) 作用。
\end{solution}

\begin{exercise}\label{b5:ex:15-induced-cyclic-relations}
设 \(\chi:\mathfrak h\to k\) 零化 \([\mathfrak h,\mathfrak h]\)，
\(k_\chi\) 为相应一维模。
证明诱导模同构于 \(U(\mathfrak g)\) 模掉由
\(x-\chi(x)\)、\(x\in\mathfrak h\)，生成的左理想。
\end{exercise}
\begin{solution}
由泛性质，\(\chi\) 唯一延拓为含幺代数同态
\(\widetilde\chi:U(\mathfrak h)\to k\)。记题中左理想为 \(J\)。
映射 \(u\mapsto u\otimes1\) 满射到诱导模，且
\(u\bigl(x-\chi(x)\bigr)\) 被送到零，故通过所述左理想商。
反过来，令 \([u]\) 表示 \(u\) 在 \(U(\mathfrak g)/J\) 中的类。
对任意 \(u\in U(\mathfrak g)\)、\(x\in\mathfrak h\)，有 \([ux]=\chi(x)[u]\)。
在 \([ux_1\cdots x_r]\) 中从最右端逐次使用此式，再按线性性延拓，得到
\[
[ua]=\widetilde\chi(a)[u]\qquad\bigl(a\in U(\mathfrak h)\bigr).
\]
因此 \(u\otimes c\mapsto c[u]\)、\(c\in k\)，满足整个 \(U(\mathfrak h)\) 上的平衡关系：
\((ua)\otimes c\) 与 \(u\otimes\widetilde\chi(a)c\) 的像均为 \(c\widetilde\chi(a)[u]\)。
两方向互逆，且保持左 \(U(\mathfrak g)\) 作用。关系只须允许左乘任意包络代数元素，不是要求它们生成双边理想。
\end{solution}

\begin{exercise}\label{b5:ex:15-verma-faithfulness}
对任意 \(\lambda\in\mathbb C\)，证明 \(\mathfrak{sl}_2\) 在 \(M(\lambda)\) 上的作用忠实，
但其包络代数作用不忠实。
\end{exercise}
\begin{solution}
Lie 代数作用的核为 \(\mathfrak{sl}_2\) 的理想。
在 PBW 基上有 \(fv_0=v_1\ne0\)，所以核不是整个 Lie 代数；
单性迫使核为零。
但中央元素 \(C-\lambda(\lambda+2)\) 作用为零，
它在包络代数中非零，因为主符号仍为 \(h^2+4fe\)。
因此包络代数作用不忠实。
\end{solution}

\begin{exercise}\label{b5:ex:15-verma-left-not-two-sided}
在 \(M(1)\) 中取最高向量 \(v\)，令
\(I=\bigl\{u\in U(\mathfrak{sl}_2)\mid uv=0\bigr\}\)。
证明 \(I\) 是左理想但不是右理想。
\end{exercise}
\begin{solution}
若 \(uv=0\)，则 \((au)v=a(uv)=0\)，故它是左理想。
\(e\in I\)，但 \(efv=hv=v\ne0\)，所以 \(ef\notin I\)。
因此 \(I\) 不封闭于右乘 \(f\)，不是右理想，更不是双边理想。
\end{solution}

\begin{exercise}\label{b5:ex:15-induction-not-simple}
给出一个正合诱导函子把简单模送到非简单模的例子，说明“保持正合”与“保持简单”不同。
\end{exercise}
\begin{solution}
取 \(\mathfrak b=\mathbb Ch\oplus\mathbb Ce\subseteq\mathfrak{sl}_2\)，
并取一维平凡 \(\mathfrak b\) 模。诱导得到 \(M(0)\)；在\cref{b5:thm:sl2-verma-submodules}中取 \(m=0\)，其由 \(v_1,v_2,\ldots\) 张成的非零真子模同构于
\(M(-2)\)，故不简单。
诱导的正合性由 PBW 右自由性保证，与其是否产生新的子模没有矛盾。
\end{solution}

\begin{exercise}\label{b5:ex:15-failed-jacobi-presentation}
设 \(k\) 为域，\(\alpha\in k\)。取含幺结合代数展示
\[
A_\alpha=k\langle x,y,z\rangle/
\bigl(xy-yx-z,\ yz-zy-x,\ zx-xz-y-\alpha z\bigr).
\]
判断哪些 \(\alpha\) 使 \(x^ay^bz^c\)、\(a,b,c\ge0\)，构成 \(A_\alpha\) 的基；求其余参数下的商代数。
\end{exercise}
\begin{solution}
在结合代数内，交换子的 Jacobi 恒等式给出
\[
[x,[y,z]]+[y,[z,x]]+[z,[x,y]]
=[x,x]+[y,y+\alpha z]+[z,z]=\alpha x=0.
\]
若 \(\alpha\ne0\)，则 \(x=0\)；第一条关系给出 \(z=0\)，第三条再给出 \(y=0\)，故 \(A_\alpha=k\)。将三个生成元都送到零的含幺同态保证单位没有被消去。
若 \(\alpha=0\)，在形式三维空间上定义交替括号 \([x,y]=z\)、\([y,z]=x\)、\([z,x]=y\)。三个互异基元的 Jacobi 和为零，含重复基元的和由交替性相消；三线性性便给出全部 Jacobi 恒等式。
于是 \(A_0\) 是这一个 Lie 代数的包络代数，由\cref{b5:thm:pbw}得到所列基。这个判断在特征二也成立，因为使用的是交替括号及其反对称性，而没有除以 \(2\)。
\end{solution}

\begin{exercise}\label{b5:ex:15-adjoint-nilpotence-degree}
设 \(x\in\mathfrak g\)，并假定 \((\operatorname{ad}x)^a=0\) 在整个 \(\mathfrak g\) 上成立。
证明对每个 \(r\ge1\)，\((\operatorname{ad}x)^{r(a-1)+1}\) 零化 \(F_rU(\mathfrak g)\)。
若 \(x^bv=0\)，给出零化 \(F_rU(\mathfrak g)v\) 的一个统一幂。
\end{exercise}
\begin{solution}
对长度至多 \(r\) 的单项式反复用导子法则。
若总作用次数为 \(r(a-1)+1\)，至少有一个因子被作用 \(a\) 次，故每项为零；
常数项也被导子零化。
再用
\(x^Nu=\sum_j\binom Nj(\operatorname{ad}x)^j(u)x^{N-j}\)。
取 \(N=r(a-1)+b\)，则每项或者 \(j\ge r(a-1)+1\)，或者 \(N-j\ge b\)，
作用于 \(v\) 都为零。这个界对整个有界滤过部分统一，但一般不能与 \(r\) 无关。
\end{solution}
